\documentclass[11pt,a4paper]{article}
\usepackage[margin=1in]{geometry}
\usepackage{amsmath,amssymb,mathtools}
\usepackage{booktabs,longtable,array,tabularx}
\usepackage{enumitem}
\usepackage{xcolor}
\usepackage{xurl}
\usepackage{microtype}
\usepackage{ragged2e}
\usepackage{fancyhdr}
\usepackage{bookmark}
\usepackage{hyperref}
\usepackage{lmodern}
\usepackage[T1]{fontenc}
\usepackage[utf8]{inputenc}
\usepackage{textcomp}
\usepackage{etoolbox}
\usepackage{hanging}

\setlength{\parindent}{0pt}
\setlength{\parskip}{6pt plus 2pt minus 1pt}
\setlength{\tabcolsep}{4pt}
\renewcommand{\arraystretch}{1.12}
\setlength{\emergencystretch}{3em}
\setlist{itemsep=2pt,topsep=4pt}
\pagestyle{fancy}
\fancyhf{}
\fancyhead[L]{DCI Pilot Preregistration}
\fancyhead[R]{v19-pilot DRAFT}
\fancyfoot[C]{\thepage}
\setlength{\headheight}{14pt}

\hypersetup{
  pdftitle={v19-pilot DRAFT: Reference-Specific, Open-Weight, Hard-Budget Directional Constraint Interference Pilot},
  colorlinks=true,
  linkcolor=blue,
  citecolor=blue,
  urlcolor=blue,
  pdfcreator={LaTeX}
}

\newcommand{\FIXED}{\textbf{FIXED}}
\newcommand{\PROV}{\textbf{PROVISIONAL}}
\newcommand{\PFD}{\textbf{PRE-FREEZE DECISION}}
\newcommand{\GATE}{\textbf{REQUIRED GATE}}
\newcommand{\SECONDARY}{\textbf{SECONDARY}}
\newcommand{\FUTURE}{\textbf{FUTURE}}
\newcommand{\NOTID}{\texttt{NOT\_IDENTIFIED}}
\newcommand{\INFERNE}{\texttt{PRIMARY\_INFERENCE\_NOT\_\allowbreak ESTIMABLE}}
\newcommand{\TECHWARN}{\nolinkurl{TECHNICAL_INTEGRITY_WARNING}}
\newcommand{\REFPOOL}{\nolinkurl{REFERENCE_POOL_INSUFFICIENT}}
\newcommand{\SATCERT}{\texttt{SAT\_CERTIFIED}}
\newcommand{\ARMMISMATCH}{\texttt{INTEGRITY\_ARM\_\allowbreak MISMATCH}}
\newcommand{\RUNFAIL}{\nolinkurl{RUN_INTEGRITY_FAILURE}}
\newcommand{\CTRLFAIL}{\nolinkurl{CONTROLLER_CAP_VIOLATION}}
\newcommand{\TERMFALSE}{\nolinkurl{TERMINAL_FALSIFICATION_CONCERN}}
\newcommand{\Rcrit}{R_{p}(b,k,o)}
\newcommand{\Tcrit}{\Theta_{\mathrm{crit}}}
\newcommand{\Bret}{\mathcal B_{\mathrm{ret}}}
\newcommand{\Sc}{\mathcal S_{\mathrm{ret}}}
\newcommand{\Jset}{\mathcal J_p}
\newcommand{\Oset}{\mathcal O}
\newcommand{\Kset}{\mathcal K}

\title{\textbf{Preregistration (Draft, Pre-Freeze)}\\[0.5em]
\large v19-pilot DRAFT: Reference-Specific, Open-Weight, Hard-Budget Directional Constraint Interference Pilot}
\author{}
\date{9 August 2026}

\begin{document}
\maketitle
\begin{abstract}
This preregistration defines a pre-freeze pilot for Directional Constraint Interference (DCI), the change in target-instruction failure probability when a focal instruction replaces one of four independently admitted matched references at fixed total load. The scientific primitive is reference-specific, $\tau_{p,j}(b,k,o)$; robustness is the least-favourable-reference functional $R_p(b,k,o)=\min_j\tau_{p,j}(b,k,o)$; and the global headline is $\Theta_{\mathrm{crit}}=\max_{s\in\mathcal S_{\mathrm{ret}}}\min_j\tau_{j,s}$ over three fixed target--focal pairs, model-specific hard reasoning budgets, paired loads $k\in\{2,5\}$, and both target/added-slot presentation orders.

The principal v19 revision is a validity-first pre-confirmatory qualification of claim-bearing finite-sample inference. The existing Maurer--Pontil empirical-Bernstein construction remains an admissible candidate, including its ordinary handling of zero empirical variance and its conservative interval completion for permanent technical missingness. A prespecified Waudby-Smith--Ramdas betting candidate is added for bounded means. Both candidates must pass theorem-assumption, missingness, multiplicity, max--min, software, and null-coverage gates under the actual DCI design before they may compete on prospective power, interval width, required superblock count, and expected compute. The selected procedure, all numerical tuning constants, multiplicity rules, implementation, and code hash are frozen before authoritative design selection and before confirmation. Stage~9 method comparison is restricted to a prospectively frozen conservative pre-Stage-10 superblock-count domain; realised Stage~10 attainability is overlaid later for transparency and can never reopen method selection. No assumed superblock count such as 150--250 is treated as established.

To make the probability model match both candidate theorems, confirmatory superblocks are generated as independent with-replacement draws from frozen pair-specific task and background populations, followed by deterministic whole-proposal feasibility rejection. The old without-replacement background cycles are removed from confirmation. The inferential probability is over this prespecified superblock-generating law and the frozen stochastic decoder; the fact that its realised pseudorandom stream is reproducible after the root seed is frozen is an implementation property, not a proof of statistical independence.

The scientific runtime is backend-neutral: PyTorch + Hugging Face Transformers + a frozen accelerator backend. ROCm remains a permitted local-development backend, while Lambda On-Demand Cloud is the planned scalable confirmatory substrate and may use a fully frozen CUDA/NVIDIA stack. The exact accelerator backend, GPU SKU, software/runtime manifest, and certified device-interchangeability class are selected by an outcome-independent correctness-and-cost benchmark before model-dependent scientific calibration. vLLM remains optional and may replace Transformers only after the complete frozen equivalence gate passes.

Permanent technical missingness remains distinct from behavioural failure and is propagated by support-valid interval completion without deleting superblocks or substituting replacement seeds. The SESOI satisfies $\epsilon\ge0.06$ from Stage~1 and its exact value is selected by a signed prospective scientific-decision procedure before any numerical procedure that depends on it. The protocol remains pre-freeze: numerical parameter distributions, admitted reference identities, checkpoint, retained budget grid, $V$, exact $\epsilon$, inference-selection outcome, runtime/hardware identities, pair-specific $H_p$, $G$, and other prespecified empirical outputs are not known until their frozen procedures are executed.
\end{abstract}
\tableofcontents
\newpage

\begin{quote}
\textbf{Protocol version:} v19-pilot DRAFT.\\
\textbf{Status:} pre-freeze preregistration. No confirmatory outcomes have been generated.\\
\textbf{Scientific unit:} one paired superblock drawn independently from the frozen pair-specific generating law; pair-specific counts $H_p$ are permitted.\\
\textbf{Primitive estimand:} one reference-specific focal-minus-reference target-failure contrast for every fixed pair, admitted reference, retained hard budget, load, and presentation order.\\
\textbf{Reference count:} exactly $J=4$ independently admitted references per target--focal pair.\\
\textbf{Compute treatment:} an externally enforced hard cap on sampled reasoning-content tokens followed by a separately bounded final-answer phase of $V$ tokens.\\
\textbf{Primary decoder:} multinomial sampling with $T=0.70$, top-$p=0.95$, top-$k$ disabled, one beam, and no additional truncation or repetition processor.\\
\textbf{Claim-bearing inference:} selected before confirmation by the frozen validity-first comparison between Maurer--Pontil empirical-Bernstein and the prespecified Waudby-Smith--Ramdas betting candidate; simultaneous primitive protection is followed by the same max--min transformation.\\
\textbf{Permanent technical missingness:} interval completion on the exact contribution support; no primary superblock deletion and no replacement seeds.\\
\textbf{Runtime:} PyTorch + Hugging Face Transformers + frozen accelerator backend; Lambda On-Demand Cloud is the planned scalable execution substrate, subject to the frozen hardware/runtime benchmark and interchangeability gate.\\
\textbf{Primary purpose:} obtain an executable, falsifiable, reference-specific DCI surface. If no feasible design satisfies the frozen prospective-power criteria, the study is declared a feasibility-and-measurement pilot without weakening the SESOI, reference panel, multiplicity protection, search space, or validity requirements.
\end{quote}

\section{Preregistration status and thesis}

\subsection{Thesis, objective, and success criterion}
Directional Constraint Interference, abbreviated \textbf{DCI}, is a causal replacement effect for a particular target--focal instruction pair under a specified model, decoder, hard reasoning budget, load, presentation order, and superblock-generating distribution. The study asks whether a focal constraint $B$ makes a target constraint $A$ more likely to fail than would a separately admitted matched reference $R_j$ occupying the same added-constraint slot:
\[
A+B+S \qquad\text{versus}\qquad A+R_j+S.
\]
The task, target, backgrounds, assigned model and runtime, hard budget, final-answer allowance, target parameter draw, canonical wording, and every other pretreatment field are paired. The intervention is the identity of the added-slot occupant together with that occupant's own frozen constraint-specific fields.

The scientific spine is fixed: three target--focal pairs, exactly four independently admitted references per pair, reference-specific primitive estimands with no primary synthetic-control averaging, least-favourable-reference robustness, a global max--min headline, paired loads $k\in\{2,5\}$, both presentation orders, model-specific hard reasoning budgets, frozen family-specific parameter distributions, and fixed-size stopping. v19 changes only those parts for which pre-freeze evidence exposed a validity, dependency-ordering, or operational defect.

The protocol succeeds operationally when two independent implementation teams given only this document and the signed freeze bundle make the same decision for every scientific draw, retry, verifier state, contribution support interval, confidence calculation, familywise multiplicity allocation, max--min transformation, result label, runtime recovery event, and terminal-falsification action.

\subsection{Protocol namespace and source precedence}
The v19 protocol occupies its own versioned namespace. Low-level provenance, retry-lineage, verifier-harness, satisfiability, schema, and causal-field components may be reused only when their frozen hashes and semantics agree with v19. Silent mutation of prior objects is prohibited.

Within v19, Section~\ref{sec:stages} controls dependency order; Section~\ref{sec:analysis} controls confirmatory inference after the method-selection stage has frozen one candidate; Section~\ref{sec:missingness} controls permanent technical missingness; and Appendix~D controls scientific-runtime qualification. Empirical-Bernstein is not the irrevocably fixed confirmatory method until the inference qualification stage resolves. Construction-only empirical-Bernstein calculations, such as the verifier-coupling gate, remain fixed where v19 does not reopen them.

No implementation may restore the discarded studentised max-$T$ bootstrap as claim-bearing inference, delete a superblock because a planned generation is permanently technically missing, replace a failed scientific seed, reintroduce a universal cross-family parameter-severity coordinate, or switch scientific runtime after confirmation begins.

\subsection{Status vocabulary}
\begin{itemize}
\item \FIXED{}: a scientific commitment or fully specified method whose alteration after freeze requires a new protocol version.
\item \PFD{}: a numerical value, artifact identity, or empirical selection output that will be produced by a frozen pre-confirmatory rule.
\item \GATE{}: a logical, statistical, measurement, software, or runtime acceptance condition that must pass before dependent work proceeds.
\item \SECONDARY{}: prespecified analysis that cannot alter the primary headline logic.
\item \FUTURE{}: outside this pilot.
\item \PROV{}: a planning candidate that has not yet passed its deterministic pre-freeze selection rule.
\end{itemize}

\subsection{Current status summary}
\begin{longtable}{p{0.43\textwidth}p{0.49\textwidth}}
\toprule
\textbf{Component} & \textbf{Current v19 status} \\
\midrule
\endhead
Replacement contrast; three fixed target--focal pairs; exactly four reference-specific primitives per pair/cell & \FIXED{} \\
Least-favourable-reference robustness and global max--min headline & \FIXED{} \\
Paired loads $k\in\{2,5\}$ and both target/added-slot orders & \FIXED{} \\
Universal cross-family parameter-severity matching coordinate & \FIXED{} excluded \\
Reference firewall, candidate pool, construction bootstrap, eligibility gates, coupling, wording, matching, ranking, and independent validation & \FIXED{} methods \\
Permanent technical missingness by interval completion; identical-draw retries; no replacement seeds; prospective complete-data conditional-mean law preserved & \FIXED{} \\
Confirmatory superblock law: independent with-replacement proposal streams plus whole-proposal feasibility rejection & \FIXED{} \\
Confirmatory inference candidate set: Maurer--Pontil empirical-Bernstein and WSR-PP betting & \FIXED{} candidate specifications; selected method is a \PFD{} \\
Inference-method validity-first selection rule, conservative pre-Stage-10 $H$ planning cap, and 10,000-experiment comparison & \FIXED{} method \\
Primitive support $[-1,1]$; compute/load/order contribution support $[-2,2]$ & \FIXED{} \\
Multiplicity guarantee & \FIXED{} separate 5\% FWER for each confirmatory family; no joint 5\% guarantee across all families \\
SESOI floor $\epsilon\ge0.06$ & \FIXED{}; exact $\epsilon$ is a \PFD{} governed by the prospective SESOI procedure and early dependency rules \\
Scientific architecture: PyTorch + Hugging Face Transformers + frozen accelerator backend & \FIXED{} \\
ROCm for local development/baselines & permitted; not required for confirmation \\
Lambda On-Demand Cloud & \FIXED{} planned scalable confirmatory substrate; exact hardware/runtime are \PFD{} / \GATE{} \\
One model process and one active scientific sequence per GPU & \FIXED{} default; tensor parallelism requires separate pre-freeze justification and validation \\
GPU interchangeability before confirmation & \GATE{}; no non-interchangeable physical-device pinning fallback \\
vLLM & optional \PROV{} optimisation; Transformers remains primary unless the full equivalence gate passes \\
Whole-composition satisfiability and concrete verified witness & \FIXED{}; exact construction artifacts are \PFD{} \\
Pair-specific $H_p$ and generation replicate count $G$ & \PFD{} from attainable-bank construction and authoritative simulation \\
Confirmatory outcomes & none generated \\
Protocol & not frozen \\
\bottomrule
\end{longtable}

\subsection{Primary purpose and prospective-power rule}
The pilot becomes confirmatory only after every construction, inference, runtime, and validation dependency has been frozen. Prospective power is evaluated by 10,000-replicate end-to-end simulation of the actual selected confidence procedure, actual family cardinalities, interval missingness, shared focal-arm dependence, superblock-generating law, and max--min headline. If no feasible $(\mathbf H,G)$ design meets the prespecified acceptance criteria, the target pairs, four-reference panels, multiplicity protection, SESOI, and search space remain unchanged. The study is frozen explicitly as a feasibility-and-measurement pilot and cannot make a powered positive-global-headline claim.

\section{The pilot in one page}

\subsection{What DCI measures}

Suppose a model is asked to perform a task while obeying a target constraint $A$. If the experiment simply compares $A+B+S$ with $A+S$, any increase in target failure mixes two causes: the particular effect of $B$ and the generic burden of receiving one more instruction. DCI therefore replaces $B$ with a matched reference $R_j$:
\[
A+B+S
\quad\text{versus}\quad
A+R_j+S.
\]
A positive contrast means that $B$ increases failure of $A$ relative to that specific admissible counterfactual.

Each reference answers a different counterfactual question. Consequently, the four reference contrasts are never averaged into a primary synthetic control. If one reference implies positive interference and another implies facilitation, the disagreement is scientific information rather than noise to be compressed away.

\subsection{Experimental axes}

The primary scientific search space contains:
\begin{itemize}
\item three fixed target--focal pairs, $p\in\{1,2,3\}$;
\item four admitted references per pair, $j\in\mathcal J_p$ with $|\mathcal J_p|=4$;
\item two paired loads, $k\in\mathcal K=\{2,5\}$;
\item a model-specific retained hard reasoning-budget set $b\in\Bret$, with $|\Bret|\ge 3$;
\item two presentation orders, $o\in\Oset=\{o_A,o_Q\}$, where $o_A$ places the target before the added-slot occupant and $o_Q$ reverses those two positions.
\end{itemize}

The confirmatory search set $\Sc$ is the complete set of pair $\times$ retained-budget $\times$ load $\times$ order regions that pass all required support, manipulation, measurement, and apparatus gates before protocol freeze. $\Sc$ is frozen before confirmatory outcomes exist.

\subsection{What one superblock contains}

A superblock contains a common pretreatment scientific context: task and topic, pair, target and focal rules, all four admitted reference identities, all sampled parameter realisations, canonical wordings, prompt template, maximum-load background set $S^{\max}$, background order, the two presentation-order renderings, and every other frozen pretreatment field.

The complete superblock is certified feasible at maximum load before either load is instantiated. The same superblock then yields
\[
k=2:\quad S=\varnothing,
\qquad
k=5:\quad S=S^{\max},\quad |S^{\max}|=3.
\]
The same task, $A$, $B$, references, parameter draws, canonical wordings, and presentation-order definitions therefore appear at both loads.

\subsection{What is measured}

The primary outcome is binary failure of the target constraint. For a delivered final answer,
\[
F_A=
\begin{cases}
0,&\text{target constraint verdict is PASS},\\
1,&\text{otherwise}.
\end{cases}
\]
A partial satisfaction of an exact target is therefore a target failure. Malformed, refused, truncated, or empty delivered outcomes are retained and do not disappear through post-treatment filtering. Technical failures that occur before a scientific output record is persisted follow the frozen retry policy. If the identical scientific draw remains unavailable, its binary verifier value is represented as $[0,1]$ and propagated through the primary contrast; the corresponding superblock remains in the analysis.

\subsection{Primitive, robust, and headline quantities}

For pair $p$, reference $j$, budget $b$, load $k$, and order $o$, the primitive effect is
\[
\tau_{p,j}(b,k,o)
=
P(F_{A_p}=1\mid A_p+B_p+S,b,k,o)
-
P(F_{A_p}=1\mid A_p+R_{p,j}+S,b,k,o).
\]

Reference robustness within one cell is
\[
R_p(b,k,o)=\min_{j\in\mathcal J_p}\tau_{p,j}(b,k,o).
\]
Thus $R_p(b,k,o)>0$ means that the focal constraint increases target failure relative to every admitted reference.

The safety-oriented headline is
\[
\Tcrit
=
\max_{(p,b,k,o)\in\Sc}
\min_{j\in\mathcal J_p}\tau_{p,j}(b,k,o).
\]
The inner minimum demands robustness to the least favourable admissible counterfactual. The outer maximum asks whether a sharp reference-robust interference regime exists anywhere in the complete preregistered search space.


\subsection{How the global confidence statement is obtained}
For every claim-bearing primitive $r$, one superblock contributes a bounded quantity $D_{ir}\in[-1,1]$. Before authoritative design selection, one finite-sample confidence procedure is selected from the two prespecified candidates in Section~\ref{sec:analysis}. The selected procedure constructs one-sided lower bounds $L_r$ simultaneously over the complete primary primitive family at familywise level $\alpha=0.05$. Permanent technical missingness enters through conservative support intervals rather than row deletion.

For a retained cell $s$, let $\mathcal B(s)$ denote its four reference-specific primitives. Define
\[
L_s=\min_{r\in\mathcal B(s)}L_r,
\qquad
L_\Theta=\max_{s\in\Sc}L_s.
\]
On the simultaneous primitive-coverage event, $L_r\le\mu_r$ for every primitive, so applying the same minimum and maximum operators preserves $L_\Theta\le\Theta_{\rm crit}$. The maximised plug-in estimate $\widehat\Theta_{\rm crit}$ is reported only as a selected descriptive estimate, directly adjacent to the claim-bearing $L_\Theta$ and the mandatory selection warning.

Zero observed variance is not an inferential failure under either candidate. The Maurer--Pontil candidate retains a positive finite-sample range term. The betting candidate uses a nonnegative capital process with smoothed predictable stakes and therefore has no division-by-observed-standard-error branch. Method selection is prospective and cannot use confirmatory DCI outcomes \href{https://arxiv.org/abs/0907.3740}{(Maurer \& Pontil, 2009)}; \href{https://doi.org/10.1093/jrsssb/qkad009}{(Waudby-Smith \& Ramdas, 2024)}.

\section{Directional Constraint Interference from first principles}

\subsection{A concrete example before notation}

Consider a target requiring exactly six sentences. The focal rule requires the sentences to increase monotonically in length. A plausible reference might impose a different, empirically matched structural burden without mechanically implying the six-sentence target. If target failure is 40\% with the focal rule and 18\% with that reference, the reference-specific DCI is $0.22$ for that condition.

Now suppose a second admitted reference yields target failure of 43\%. The corresponding DCI is $-0.03$. Averaging these references would produce $0.095$, a number that hides the counterfactual dependence. The v19 protocol instead reports both effects. A reference-robust positive claim fails because the least favourable reference does not support positive interference.

\subsection{Why replacement identifies a different question from load evaluation}

Ordinary constraint-load evaluation asks how performance changes as more requirements are imposed. That quantity is useful, but it does not identify whether a particular constraint $B$ causes failure of $A$ beyond generic additional-instruction burden. Replacement holds total load fixed. The causal contrast concerns the identity of the added-slot constraint rather than the presence versus absence of a slot.

This distinction also separates DCI from symmetric co-failure statistics. A symmetric association between failures of $A$ and $B$ can be large without $B$ causing $A$ to fail, and a directional causal effect can exist even when aggregate co-failure association is small. The motivating Constraint Saturation Evaluation literature therefore supplies background evidence about compositional difficulty, not an estimator of the DCI effect itself \href{https://openreview.net/forum?id=a6r2ECEVls}{(Vasileva, 2026)}.

\subsection{Why each reference is scientifically meaningful}

An admitted reference $R_j$ represents a separately validated counterfactual occupant of the same instructional slot. The comparison with $R_1$ is not interchangeable with the comparison with $R_2$ because the two references can differ in residual characteristics even after matching. Reference construction is designed to bound those differences, not to prove that all references are identical.

The four-dimensional causal profile
\[
\boldsymbol\tau_p(b,k,o)
=
\left(
\tau_{p,1}(b,k,o),\ldots,\tau_{p,4}(b,k,o)
\right)
\]
is therefore the fundamental result for a cell. Generation replicates may be averaged because they repeatedly measure the same causal condition. Scientifically distinct references, budgets, loads, orders, and pairs are not averaged merely to obtain a simpler scalar.

\subsection{Why the least favourable reference is the robustness operator}

For a given cell, the minimum
\[
R_p(b,k,o)=\min_j\tau_{p,j}(b,k,o)
\]
answers a strong question: how large is the focal effect even under the admitted counterfactual that makes the case for interference weakest? A positive value cannot be produced by choosing a convenient reference after seeing outcomes.

The minimum is not tested as an independent second hypothesis. Its uncertainty is derived deterministically from simultaneous confidence bounds on the primitive effects. This prevents double-counting multiplicity and keeps the inferential primitive aligned with the causal primitive.


\paragraph{Worked example: least-favourable reference.}
Suppose a cell has four estimated primitive effects $(0.18,0.12,0.25,0.07)$ and simultaneous lower bounds $(0.08,0.01,0.13,-0.02)$. The point-estimate robustness is $\\widehat R=0.07$, yet the robust lower bound is $L^R=-0.02$. The cell therefore does not support reference-robust positive DCI. No average of the four estimates can overturn the least-favourable simultaneous bound.

\subsection{Why the critical-regime maximum is protected globally}

The outer maximum in $\Tcrit$ searches the preregistered grid for the sharpest reference-robust regime. A maximum preferentially selects upward sampling noise if treated naively. The protocol therefore protects all primitive effects in the complete frozen search space simultaneously. Only after those bounds are constructed is the deterministic maximum of least-favourable lower bounds calculated.

This is a prespecified search functional, not outcome-dependent design selection. The pair, budget, load, order, and reference sets are frozen before confirmation. The study reports the full causal surface in addition to the maximising condition so that the headline cannot substitute for the evidence it summarises.

\section{Experimental design}

\subsection{Fixed factorial structure}

The confirmatory study uses one selected open-weight checkpoint. Within that checkpoint, every frozen superblock is evaluated at every retained hard reasoning budget, both loads, both presentation orders, the focal arm, and all four reference arms. Replicates repeat stochastic generation within an otherwise identical cell.

The design is therefore crossed within superblock over budget, load, presentation order, and arm. Pair is a fixed between-superblock stratum because each superblock belongs to one of the three prespecified target--focal pairs.

\subsection{Presentation order is restored as a primary moderator}

\FIXED: presentation order is a crossed primary moderator. Two order levels are evaluated for every superblock:
\begin{itemize}
\item $o_A$: the target constraint $A$ appears before the occupant of the added slot;
\item $o_Q$: the added-slot occupant $Q\in\{B,R_1,\ldots,R_4\}$ appears before $A$.
\end{itemize}

Background constraints retain the same frozen relative order and prompt positions under both order renderings. The two focal constraints swap only the target and added-slot positions according to the template's explicit order rule. No primary analysis averages $o_A$ and $o_Q$ before inference. If the effect reverses sign by order, the reversal remains visible in the primitive surface and in the critical-regime search.

This choice increases the multiplicity family and generation count. That cost is accepted because an order average can turn two large opposing causal effects into a misleading near-zero effect.

\subsection{Pair dependence is descriptive, not population-level}

The three pairs are purposively selected scientific cases. They are not exchangeable random draws from a population of constraint pairs. Pair identity is perfectly confounded with the selected target/focal family combination. The pilot may compare their reference-specific surfaces descriptively. It may not infer a population-level family effect or claim that relational constraints as a class have a different DCI distribution from lexical or structural constraints.

A future design with multiple independently sampled pairs per family combination is required for family-level population inference.

\subsection{Budget is categorical and model-specific}

The hard reasoning budget is a categorical treatment. No monotonicity assumption is imposed. The primary compute contrast compares the highest and lowest retained budgets, while every retained budget remains visible in the primitive surface. Adjacent-budget contrasts are prespecified secondary summaries.

Raw reasoning-token counts are interpreted only within the selected model and tokenizer. A token under one tokenizer is not treated as the same cardinal compute unit as a token under another model family. Cross-model scale-up is therefore replication under different model-specific budget grids, not a pooled dose-response analysis.

\subsection{Load is paired inside superblock}

Load dependence is a within-superblock contrast. The maximum-load environment is constructed first, then the same scientific context is instantiated at $k=2$ by removing the three background constraints and at $k=5$ by restoring them. This design eliminates the previous comparison between independently filtered low-load and high-load block populations.

The interpretation remains conditional on the maximum-load-feasible population. The study does not claim transport to arbitrary tasks that could never support the $k=5$ composition.


\paragraph{Worked example: paired load dependence.}
For one pair/reference/order/budget, suppose $\\widehat\\tau(b,2)=0.04$ and $\\widehat\\tau(b,5)=0.19$. The load contrast is $0.15$. Its uncertainty is computed from the within-superblock difference because the same task, target, focal/reference identities, parameter draws, and canonical wording generate both loads. Treating the two loads as independent samples would discard this pairing and answer a different question.

\paragraph{Worked example: compute can change DCI without monotonically changing performance.}
Suppose target failure in the focal and one reference arm changes from $(0.35,0.30)$ at $b_{\\min}$ to $(0.20,0.05)$ at $b_{\\max}$. Both arms improve in absolute terms, but reference-specific DCI rises from $0.05$ to $0.15$, so $\\Delta_b\\tau=0.10$. Better task performance therefore does not imply reduced interference.

\section{The three target--focal pairs}

\subsection{Rule definitions}

\begin{longtable}{p{0.06\textwidth}p{0.27\textwidth}p{0.27\textwidth}p{0.30\textwidth}}
\toprule
\textbf{Pair} & \textbf{Target $A$} & \textbf{Focal $B$} & \textbf{Descriptive family combination} \\
\midrule
\endhead
1 & Produce exactly $N$ bullet-list items. & Place the items in the correct transitive order. & relational/order $\rightarrow$ format \\
2 & Include all specified mandatory words. & Every word in the response must be unique. & word-level/global uniqueness $\rightarrow$ lexical inclusion \\
3 & Produce exactly $N$ sentences. & Sentence lengths must increase monotonically. & ordering $\rightarrow$ structural sentence count \\
\bottomrule
\end{longtable}

The rule definitions are \FIXED. Existing registry identifiers are not treated as authoritative where they collide or are unverified. A pre-freeze identifier reconciliation must assign unique v19-native constraint identifiers without changing the rule semantics.

\subsection{Pair 1: transitive ordering \texorpdfstring{$\rightarrow$}{to} exact bullet count}

The target requires exactly $N$ bullet items. The focal rule imposes a specified transitive ordering relation among the items. The principal measurement risk is verifier coupling because both the count and ordering verifiers may parse the same list representation. A malformed list could therefore change both verdicts through shared parsing rather than through genuine interference.

Pair 1 requires a deterministic target verifier, deterministic focal verifier, canonical and held-out semantic wordings, a scope-compatible candidate reference pool, unrelated-target banks, a pair-specific coupling probe bank, calibrated parameter distributions, and human reliability audits before it can enter confirmation.

\subsection{Pair 2: all-words-unique \texorpdfstring{$\rightarrow$}{to} mandatory-word inclusion}

The target requires a sampled set of mandatory words to appear. The focal rule requires every word in the final answer to be unique under the frozen normalisation and tokenisation-for-verification rule. The principal scientific risks are logical incompatibility and floor effects. A parameter realisation that forces repeated mandatory tokens is outside the DCI estimand because the composition is impossible. A focal parameter regime that is almost never satisfiable is also uninformative as a manipulation.

Parameter calibration and whole-composition satisfiability therefore operate on actual sampled mandatory-word sets and uniqueness semantics, not only on the constraint-family labels.

\subsection{Pair 3: increasing sentence lengths \texorpdfstring{$\rightarrow$}{to} exact sentence count}

The target requires exactly $N$ sentences. The focal rule requires successive sentence lengths to increase monotonically under the frozen sentence segmentation and length unit. Existing infrastructure from the preceding experiment may be reused only if its semantics and hashes are unchanged and it passes the v19 reliability, coupling, wording, parameter-distribution, reference-certificate, and superblock requirements.

The principal measurement risk is shared sentence segmentation. Sentence-boundary errors and conversational preambles can affect both target and focal verdicts. Pair-specific coupling audits remain mandatory.

\subsection{What the pair labels cannot establish}

The family labels are descriptive metadata. Because each family combination is represented by one selected pair, cross-pair differences can arise from the exact constraints, parameter distributions, task scopes, wordings, verifiers, reference panels, or background pools. The pilot therefore does not test a scalar pair-heterogeneity hypothesis and does not infer family-level effects.

\section{Parameterised constraints}

\subsection{Parameters vary under frozen distributions}
For every parameterised constraint $Q$, define an admissible finite or continuous set $\mathcal P_Q$ and a sampling distribution $\nu_Q$,
\[
\theta_Q\sim\nu_Q,\qquad \theta_Q\in\mathcal P_Q.
\]
The parameter spaces are family-specific. No scalar ``severity'' coordinate is constructed to make heterogeneous quantities such as sentence count, mandatory-word count, ordering depth, or whole-response uniqueness commensurable.

\subsection{Boundary-aware parameter calibration}
Before any confirmatory generation, each parameterised family receives a predeclared ordered grid of at most five admissible settings and an explicit easier/harder direction. The canonical setting is tested first. Each required calibration cell uses $n_{\rm cal}=100$ independent observations and the exact two-sided 95\% Clopper--Pearson interval. A cell passes only when
\[
CI^{CP}_{0.95}(p)\subseteq[0.05,0.95].
\]
For $n=100$, 11 through 89 successes inclusive satisfy this rule. Exact binomial procedures are used because normal approximations behave poorly near boundaries \href{https://www.stat.purdue.edu/~dasgupta/publications/tr99-19.pdf}{(Brown, Cai, \& DasGupta, 2001)}.

If a cell is too difficult, move exactly one step in the frozen easier direction. If it is too easy, move exactly one step in the frozen harder direction. The first passing value is selected. Interpolation, manual parameter tuning, and post hoc choice are prohibited. If no value passes, a candidate reference is hard-ineligible; a required target/focal item is replaced only by a prespecified reserve item before freeze; and an unsalvageable proposed composition is excluded from the construction pool.

The screen applies to each target, focal, varying background component, hard-eligible reference, and every instantiated $A+B+S$ and $A+R_j+S$ regime under every budget, load, and order capable of entering $\Sc$. Confirmatory floor or ceiling outcomes never trigger retrospective recalibration or exclusion.

\subsection{Frozen parameter population}
Only after the boundary gate and satisfiability requirements have been applied is a nonempty $\mathcal P_Q$ and explicit $\nu_Q$ frozen. The confirmatory estimand therefore averages over the \emph{pre-frozen admissible population} induced by these construction rules. This restriction is part of the target population, not a post-treatment filter.

\subsection{Pairing of realised values}
Within a superblock, each constraint identity receives one realised parameter draw. The target draw is identical across focal/reference arms, budgets, loads, orders, and replicates. The focal draw is held fixed wherever the focal appears; each reference has its own fixed draw; and each background constraint has one draw used in all $k=5$ arms. Parameter randomisation therefore contributes to the superblock-generating distribution while avoiding unnecessary within-superblock noise.

\subsection{Parameter severity is excluded from matching}
\begin{quote}\textbf{Parameter severity as a cross-family reference-matching coordinate: FIXED, excluded.}\end{quote}
Reference comparability is established using the empirical matching coordinates in Section~\ref{sec:references}. Family-specific parameters and descriptors are still archived and reported. Matching summaries are marginalised over each constraint's own frozen $\nu_Q$, while whole-composition satisfiability and boundary eligibility are checked on realised compositions.

\section{Reference constraints and counterfactual validity}\label{sec:references}

\subsection{Why references are necessary}
Deletion changes total load. DCI therefore compares the focal constraint with matched alternative constraints occupying the same slot. Each admitted reference is a distinct counterfactual, so primary inference never averages references into a synthetic control.

\subsection{Candidate-pool rule}
For each target--focal pair, author 24 candidates under the frozen rubric. Apply every hard logical, measurement, satisfiability, boundary, coupling, and semantic-wording gate before statistical matching. At least 12 candidates must survive all hard gates. If fewer than 12 survive, author exactly one additional blinded tranche of 12 and apply the identical gates. If fewer than 12 of the resulting maximum of 36 survive, assign \REFPOOL{} and do not freeze the affected pair. No gate, caliper, certificate, or verifier criterion is relaxed to achieve the minimum.

\subsection{Matching coordinates}
Reference comparison uses four empirical properties, measured on unrelated targets and marginalised over the candidate's frozen parameter distribution: isolated difficulty $d$; outgoing burden $g_{\rm burden}$, the disruption caused to unrelated targets; incoming vulnerability $r_{\rm in}$, the candidate's own degradation under unrelated added constraints; and wording instability $W$, defined below. Instruction length is reported and is not matched because literal canonical wording is part of the instruction-level treatment. Generative scope, output-schema compatibility, satisfiability, verifier validity, and coupling are hard eligibility conditions.

\subsection{Construction firewall}
Reference authoring, hard eligibility, calibration, matching, ranking, and held-out validation may use only dedicated pre-confirmatory material. They may not inspect confirmatory outcomes, realised reasoning use, answer length, truncation, refusal, malformed-output rates, co-failure correlations, DCI estimates, or candidate-pool spread as a justification for changing tolerances. Local baseline DCI effects are also sequestered from checkpoint, SESOI, reference, $H_p$, and $G$ selection.

\subsection{Construction bootstrap unit}
Define one independent construction cluster $U_i$ as one base calibration task/composition together with all parameter realisations, candidate references, wording/paraphrase observations, matching coordinates, relevant budgets, presentation orders, and generation replicates attached to that base unit. All of those nested observations travel together whenever $U_i$ is resampled.
Use exactly $B=10{,}000$ bootstrap replicates. Each bootstrap draws complete $U_i$ units with replacement and recomputes every candidate's matching vector and admission statistic from scratch. Individual generations, candidate rows, budgets, orders, parameter realisations attached to a task, and nested paraphrases are never resampled independently. This preserves the dependence structure created by the construction sampling design, consistent with cluster-resampling principles \href{https://cameron.econ.ucdavis.edu/research/Cameron_Miller_JHR_2015.pdf}{(Cameron \& Miller, 2015)}. The bootstrap PRNG, seed, software version, and resampling code hash are frozen.

\subsection{Simultaneous uncertainty-aware coordinate calipers}
For candidate $C$, coordinate $r$, budget $b$, and order $o$, define
\[
\Delta_{C,r,b,o}=z_{B,r}(b,o)-z_{C,r}(b,o).
\]
Construct simultaneous bounds $[L_{C,r,b,o},U_{C,r,b,o}]$ over every candidate actually searched, every coordinate, calibration budget, and primary order. Define
\[
U^{\rm abs}_{C,r,b,o}=\max\{|L_{C,r,b,o}|,|U_{C,r,b,o}|\}.
\]
Candidate $C$ passes coordinate $r$ only if
\[
\boxed{\max_{b,o}\frac{U^{\rm abs}_{C,r,b,o}}{\delta_r(b,o)}\le1.}
\]
The numerical calipers $\delta_r$, sensitivity coefficients, and reference-error allowance remain pre-freeze numerical values linked to the SESOI and calibration evidence; they are not fitted to force four candidates to pass.

\subsection{Joint reference-error certificate}
For frozen matching families $f$,
\[
\operatorname{Cert}_{b,o}(C)=\sum_f s_f(b,o)U^{\rm abs}_{C,f,b,o}.
\]
Each family is charged once. Admission additionally requires
\[
\boxed{\max_{b,o}\frac{\operatorname{Cert}_{b,o}(C)}{\varepsilon_{\rm ref}(b,o)}\le1.}
\]
Budgets and orders are alternative conditions, so the certificate takes a worst-case maximum rather than adding across them.

\subsection{Verifier-coupling audit}
Deterministic rule-based evaluation is useful for instruction following, as illustrated by IFEval and ComplexBench \href{https://arxiv.org/abs/2311.07911}{(Zhou et al., 2023)}; \href{https://arxiv.org/abs/2407.03978}{(Wen et al., 2024)}. Shared parsers or formatting assumptions can nevertheless induce correlated measurement error.

For target $A$ and candidate $R$, define known-ground-truth verifier-error indicators
\[
E_A=\mathbf1\{\widehat V_A\ne Y_A\},\qquad E_R=\mathbf1\{\widehat V_R\ne Y_R\}.
\]
For perturbation class $t$ define
\[
\kappa_t=\left|P(E_A=1,E_R=1)-P(E_A=1)P(E_R=1)\right|.
\]
Before candidate audits, author 2,000 independent base outputs approximately balanced across $(Y_A,Y_R)\in\{(0,0),(0,1),(1,0),(1,1)\}$. Apply eight frozen truth-preserving nuisance classes: whitespace/blank lines; Markdown wrappers; benign code fences; Unicode normalization; irrelevant punctuation; irrelevant capitalization; list-marker/enumeration variation; and field ordering/whitespace for order-insensitive structured objects. Inapplicable transformations are replaced by prevalidated examples before audit results are viewed. The same probe bank is used for every candidate.

Construct one simultaneous bounded-outcome empirical-Bernstein confidence system across every audited candidate, all eight perturbation classes, and each target-error $p_{A,t}$, reference-error $p_{R,t}$, and joint-error $q_t$ probability. Given rectangles for these probabilities, evaluate
\[
U_{\kappa,t}=\max_{p_A,p_R,q}|q-p_Ap_R|
\]
exactly over rectangle boundary combinations. A candidate passes only if
\[
\max_t U_{A,t}\le0.02,\qquad \max_t U_{R,t}\le0.02,
\]
and
\[
\boxed{\max_tU_{\kappa,t}\le\min(0.015,\epsilon/4).}
\]
If 2,000 probes cannot certify the SESOI-linked threshold, increase the base-probe count to the smallest multiple of 500 for which the zero-error simultaneous empirical-Bernstein UCB is no greater than $\epsilon/8$. The threshold never changes. This gate controls correlated measurement error; it does not redefine DCI as a coupling-adjusted effect.

\subsection{Semantic paraphrase certification and wording instability}
For every target, focal, and surviving reference constraint, freeze $M=12$ semantically equivalent wordings: one canonical wording and 11 paraphrases. Paraphrases are created before wording-evaluation outcomes are inspected and must preserve numerical values, quantifiers, scope, negation, modality, precedence, and structured-output semantics. Two independent human reviewers must agree on semantic equivalence; disagreement rejects the paraphrase before evaluation.

Let $z_m=E[Y\mid w_m]$ be the satisfaction probability for wording $m$. The population wording instability is
\[
W=\sqrt{\frac{1}{M-1}\sum_{m=1}^M(z_m-\bar z)^2}.
\]
Evaluate each wording on the same $N_W=100$ independent calibration tasks with $G_W=2$ paired stochastic replicates, using common task identities and paired replicate indices across wordings. Small semantic-preserving prompt changes can alter LLM performance, motivating an explicit wording-robustness coordinate \href{https://arxiv.org/abs/2406.19898}{(Wahle et al., 2024)}; \href{https://arxiv.org/abs/2306.04528}{(Zhu et al., 2023)}.

Let $\widehat{\mathbf z}=(\hat z_1,\ldots,\hat z_M)^\top$, $P=I-M^{-1}\mathbf1\mathbf1^\top$, and
\[
S_z^2=\frac{1}{M-1}\widehat{\mathbf z}^{\top}P\widehat{\mathbf z}.
\]
Bootstrap the 100 task units as clusters, carrying all 12 wordings and both replicates together. With bootstrap covariance $\widehat\Sigma$,
\[
\widehat W=\sqrt{\max\left\{0,S_z^2-\frac{1}{M-1}\operatorname{tr}(P\widehat\Sigma)\right\}}.
\]
Use 10,000 task-cluster replicates to obtain the simultaneous one-sided UCB for $|W_R-W_A|$. The wording coordinate passes only when
\[
\boxed{U_{|W_R-W_A|}\le\min(0.02,\epsilon/2).}
\]

\subsection{Ranking, simultaneous admission, and independent validation}
After hard eligibility, coupling, wording, coordinate calipers, and the joint certificate, rank only already admissible candidates by the deterministic caliper-normalised distance
\[
d_C=\left[\sum_{r\in\{d,g_{\rm burden},r_{\rm in},W\}}
\left(\max_{b,o}\frac{U^{\rm abs}_{C,r,b,o}}{\delta_r(b,o)}\right)^2\right]^{1/2}.
\]
Select the four smallest $d_C$. This ranking cannot rescue a candidate that failed any gate or certificate and introduces no new tolerance. Exact ties use the frozen candidate-ID lexicographic/hash rule. The complete eligible set, distances, four selected references, and rejection reason codes are exported.

A disjoint validation split then repeats logical eligibility, support, verifier/coupling, wording, and matching-support checks under simultaneous uncertainty. One reference's excess error cannot be offset by another reference. Failure of any selected reference returns the pair to pre-freeze construction under the already frozen candidate-pool process; no confirmatory outcome exists at that point.

\section{Background constraints, satisfiability, and paired superblocks}

\subsection{Compositional loads}
The compositional load $k$ is the number of simultaneously active constraints. The pilot uses
\[
\mathcal K=\{2,5\}.
\]
At $k=2$, only the target and added-slot occupant are active. At $k=5$, the same target and occupant are accompanied by exactly three background constraints. Maximum-load feasibility is established first; the $k=2$ member is the paired restriction obtained by removing those backgrounds.

\subsection{Frozen background and task populations}
Before confirmatory superblocks are drawn, each pair $p$ receives a frozen confirmation task population $\mathcal T_p$ and a scope-compatible background pool $\mathcal G_p$. The background pool is partitioned into exactly three strata: lexical/content, structural/format, and relational/order. Each entry records a unique constraint ID, stratum, generative scope, output-schema requirements, verifier/audit hashes, canonical wording, parameter distribution if applicable, incompatibility and entailment metadata, and construction provenance. A constraint belongs to only one background-sampling stratum.

Neither pool may contain the pair's target, focal constraint, admitted references, semantic duplicates, known incompatible rules, or content allocated to another scientific split. The content-sufficiency gate must establish enough unique material for calibration, reference construction, validation, confirmation, and terminal falsification before expensive model-dependent work begins. Repeated confirmatory draws from a frozen finite population are allowed because the inferential population is the probability distribution over draws, not a requirement that every realised task identity appear at most once.

\subsection{Independent with-replacement superblock proposal law}\label{sec:bgsampling}
\FIXED{} v19 uses independent proposal streams across superblocks. This replaces deterministic without-replacement background cycles because the claim-bearing confidence candidates require an explicitly stated probability model.

For each pair $p$ and intended superblock index $i$, generate a fresh proposal independently as follows:
\begin{enumerate}
\item Draw one task/topic $T_{pi}$ uniformly with replacement from the frozen $\mathcal T_p$, unless Stage~1 freezes a different explicit probability mass function over $\mathcal T_p$ before any model-dependent calibration. The same probability mass function is used for every $i$.
\item Independently draw exactly one eligible background identity with replacement from each of the three strata of $\mathcal G_p$, using the frozen within-stratum probability mass functions. Uniform within stratum is the default and any departure must be frozen at Stage~1.
\item Draw each active constraint's parameter independently from its own frozen $\nu_Q$, subject to the within-superblock pairing rules in Section~6.4.
\item Draw the background-position permutation and any other stochastic pretreatment construction field independently from its frozen distribution.
\item Evaluate the complete proposed maximum-load composition in the focal arm and all four reference arms. If any required composition fails the prespecified whole-composition satisfiability or construction boundary gate, reject the entire proposal and independently draw a new proposal under the same law. No individual background identity is swapped in isolation.
\item The first proposal that passes every required gate becomes accepted superblock $i$.
\end{enumerate}

This is ordinary rejection sampling from a fixed proposal law. Consequently, accepted superblocks are independent draws from the proposal distribution conditional on the frozen pre-confirmatory acceptance event, provided the per-index proposal streams are independent. The implementation uses domain-separated, counter-based SHA-256 seed derivation for each $(p,i,\text{proposal attempt},\text{field})$ so that retrying or parallelising one superblock cannot consume random draws assigned to another. The mathematical assumption is independence under the stated design distribution. Reproducibility of the realised pseudorandom stream after a root seed is frozen does not itself establish statistical independence.

The root construction seed and canonical sampler implementation are frozen before the first confirmatory superblock is realised. A complete proposal ledger records every accepted and rejected proposal. Confirmatory outcome magnitudes cannot alter sampling probabilities, acceptance rules, or the number of proposal attempts.

\subsection{What is random under the inferential model}
For a fixed pair, the superblock contribution is random through the independent task draw, background identities, parameter draws, construction permutations, and independent stochastic decoder draws attached to the resulting superblock. All scientific rules, population supports, probability mass functions, reference identities, checkpoint, runtime, decoder, and verifiers are fixed before confirmation. The estimand averages over this frozen pair-specific data-generating law.

Conditional on past accepted superblocks, the next accepted superblock has the same frozen distribution. Therefore, for any fixed primitive $r$ within pair $p$,
\[
E[D_{ir}\mid\mathcal F_{i-1}]=\tau_r,
\]
and the stronger independent-and-identically-distributed condition required by the Maurer--Pontil candidate is also satisfied by design. Generation replicates $g=1,\ldots,G$ are averaged inside $D_{ir}$ and are not treated as independent superblocks. The Waudby-Smith--Ramdas betting construction uses the conditional-mean property directly \href{https://doi.org/10.1093/jrsssb/qkad009}{(Waudby-Smith \& Ramdas, 2024)}.

Waudby-Smith and Ramdas also derive betting procedures for uniform sampling without replacement from a fixed finite population. v19 does not invoke that result for confirmation because the DCI superblock includes stochastic task, parameter, feasibility, and decoder components rather than merely revealing a fixed finite list in random order. The confirmatory design instead removes the mismatch prospectively by using the with-replacement generating law above.

\subsection{How satisfiability affects a proposed draw}
A proposed maximum-load background triplet is accepted only when the complete superblock is feasible for $A+B+S$ and every $A+R_j+S$ arm at the actual sampled parameters. If any required arm is impossible, uncertain, times out under a required solver check, lacks required witness evidence, or fails a frozen construction boundary gate, the entire proposal is rejected. The next proposal is drawn from scratch under the same frozen law. Every proposal, rejection reason, solver record, witness, verifier output, and attempt seed is archived.

If proposal acceptance becomes so rare that the attainable-bank requirements cannot be met under the frozen design, construction fails before confirmation. The remedy is additional disjoint content or a new protocol version, not relaxation of eligibility, matching, satisfiability, SESOI, multiplicity, or inference requirements.

\subsection{Whole-composition satisfiability and evidential hierarchy}
Every prospective $A+B+S$ and $A+R_j+S$ composition must receive \SATCERT{} before entering a frozen confirmatory bank. Individual constraint validity and pairwise compatibility are insufficient.

For machine-expressible clauses, encode the conjunction in Z3 or an equivalently frozen solver. Z3 distinguishes \texttt{sat}, \texttt{unsat}, and \texttt{unknown}; only \texttt{sat} may support progression, while \texttt{unsat}, \texttt{unknown}, timeout, or missing required solver output fails closed \href{https://z3prover.github.io/api/html/js/interfaces/Solver.html}{(Z3 Project, 2026)}. A \texttt{sat} result establishes satisfiability only of the encoded abstract feature projection. It does not establish that a natural-language answer satisfying the full intended semantics exists.

The stronger operational evidence is a concrete complete witness $x^\ast$ for the actual natural-language composition. Witnesses may be generated automatically. The permitted pipeline is: (1) generate a candidate witness by a deterministic constructor or a separately seeded pre-confirmatory generator; (2) run the complete deterministic verifier conjunction on that same witness; (3) retain the witness only if every active deterministically verifiable constraint passes. Human semantic review is required only for clauses whose intended semantics are not exhaustively covered by the frozen machine encoding and deterministic verifier. Such clauses receive two independent reviews, with a third reviewer adjudicating disagreement; unresolved disagreement fails certification.

The evidential hierarchy is therefore: solver \texttt{sat} for the encoded abstraction; concrete all-verifier-pass witness for operational natural-language satisfiability; and human semantic review where deterministic coverage is incomplete. No weaker item substitutes for a stronger item when the latter is required. Archive the exact composition, structured encoding, solver input/output, candidate/winning witness, verifier results, semantic reviews where applicable, hashes, and construction commit.

Each constraint template also receives at least 1,000 automated boundary mutations covering applicable off-by-one, duplication, omission, forbidden-item, ordering, formatting, Unicode, and parser-edge cases. Property-based testing is evidence about implementation invariants rather than a proof of natural-language semantics \href{https://doi.org/10.1145/3360607}{(Lampropoulos, Hicks, \& Pierce, 2019)}.

\subsection{Superblock definition and repeated contexts}
Accepted superblock $i$ for pair $p$ contains one accepted draw of: task/topic; $A_p$, $B_p$, and the four frozen references; one parameter realisation for every active identity; canonical wordings; prompt template; maximum-load background set $S_i^{\max}$; background-position permutation; both target/added-slot order renderings; and every other frozen pretreatment field.

Because sampling is with replacement, two accepted indices may legitimately draw identical task/background/parameter content. Define a \texttt{context\_fingerprint} as the hash of those pretreatment contents. Define the scientific unit identifier separately as
\[
\texttt{superblock\_id}=\operatorname{hash}(\texttt{protocol},p,i,\texttt{context\_fingerprint}).
\]
Thus repeated contexts are distinct independent draws rather than accidental identifier collisions. The accepted draw index $i$ is never reused or renumbered after the bank is frozen.

Each superblock is crossed over both loads, every retained budget, both presentation orders, the focal plus four reference arms, and $G$ decoder replicates. The claim-bearing contribution is formed at superblock level. A permanently unavailable constituent generation changes that contribution into an interval but does not remove the superblock.

\subsection{The lower load has one paired instantiation}
At $k=2$, $S=\varnothing$. There is exactly one lower-load instantiation paired with each accepted maximum-load superblock. The task, target, focal/reference identities, parameters, canonical wordings, and order definitions are therefore common across the two load conditions.

\subsection{Arm identity is an execution invariant}
Focal and reference arms are required to be identical on every shared pretreatment field except the added-slot occupant and that occupant's own frozen constraint-specific parameter/wording fields. Any discrepancy in a field that should be common emits \ARMMISMATCH{} and is an integrity failure. No balance statistic is used to excuse such a mismatch.

Population diagnostics report task/background occupancy, duplicate-context rates, proposal rejection rates, and accepted-versus-rejected composition summaries. These are descriptive diagnostics of the frozen rejection-sampling law and never become post hoc reweighting variables.

\section{Open-weight model ladder, hard reasoning budgets, and scientific runtime}

\subsection{Design principle}
Confirmation uses one pinned open-weight checkpoint whose weights, tokenizer, configuration, chat template, and immutable revision can be archived and hashed. Small open-weight checkpoints from multiple families may be used for local infrastructure baselines. The confirmatory checkpoint may use remote compute. Baseline DCI outcomes are sequestered from checkpoint, reference, SESOI, inference-method, $H_p$, $G$, and hardware selection.

The scientific software architecture is
\[
\boxed{\text{PyTorch + Hugging Face Transformers + frozen accelerator backend}.}
\]
The accelerator backend is part of the implemented stochastic intervention because numerical kernels can affect logits and therefore sampled tokens. It must be frozen before model-dependent scientific calibration and cannot change during confirmation.

\subsection{Local development and planned confirmatory substrate}
ROCm is permitted for local development and open-weight baselines on compatible AMD hardware. It is not a requirement for confirmation. Lambda On-Demand Cloud is the planned scalable confirmatory execution substrate. A CUDA/NVIDIA backend is therefore explicitly permitted if it passes the frozen pre-confirmatory qualification and hardware-selection procedure.

The default confirmatory architecture is one model process and one active scientific sequence per GPU. Parallelism is obtained by assigning distinct immutable generation IDs to independent GPU workers. This avoids making ordinary tensor parallelism or serving-engine batching a prerequisite for feasibility. A checkpoint that fits the complete maximum prompt plus $b_{\max}+V$ plus control-token envelope on one eligible GPU is preferred. Tensor parallelism is prohibited unless, before model-dependent calibration, a separate written justification establishes that no scientifically acceptable single-GPU checkpoint is available and the complete cross-device, seed, controller, failure-recovery, and equivalence suite passes under the exact tensor-parallel configuration.

Lambda implementation facts are volatile and must be reverified immediately before freeze. The protocol relies only on capabilities established by the selected runtime benchmark, not on a preregistered current SKU, throughput, availability, or price. Lambda documents on-demand instances and regional persistent filesystems as separate cloud resources; the exact service configuration used by confirmation is recorded in the freeze manifest \href{https://docs.lambda.ai/public-cloud/on-demand/}{(Lambda, 2026a)}; \href{https://docs.lambda.ai/public-cloud/filesystems/}{(Lambda, 2026b)}.

\subsection{Checkpoint eligibility}
A checkpoint is confirmatory-eligible only if, before model-dependent scientific calibration: its weights/configuration are downloadable and hashable under the applicable licence; exact model/tokenizer/template revisions are pinned; reasoning and final-answer segments are deterministically separable under the frozen adapter; the hard cap can be enforced without candidate-answer selection; $V$ is independent of assigned $b$; at least three distinguishable retained hard-budget conditions can survive the manipulation gate; the complete runtime/dtype path passes Stage~2A correctness tests; and the maximum frozen prompt plus control tokens, $b_{\max}$, and $V$ fits the candidate hardware without OOM.

Checkpoint selection uses only predeclared capability, architecture-diversity, single-device-fit, licensing, reproducibility, and resource criteria. Observed confirmatory or baseline DCI magnitude cannot select the checkpoint.

\subsection{Outcome-independent Lambda hardware selection}\label{sec:hardwareselection}
After a checkpoint/controller/runtime candidate has passed the early scientific-runtime gates, benchmark every Lambda On-Demand GPU SKU in the frozen candidate list that can plausibly fit the model. The candidate list is generated immediately before benchmarking from Lambda's current instance-type API, which reports offered instance types, specifications, pricing, and regional capacity, and is archived with retrieval date and raw response \href{https://docs.lambda.ai/public-cloud/cloud-api/}{(Lambda, 2026c)}. Do not hard-code H100, A100, B200, GPU count, current price, throughput, or wall-clock duration in the scientific protocol.

Each candidate configuration must satisfy all of the following before cost ranking:
\begin{enumerate}
\item the complete maximum prompt plus control-token, $b_{\max}$, and $V$ envelope fits with the frozen dtype and attention backend;
\item the complete Stage~2A correctness, determinism, replay, persistence, and recovery suite passes;
\item no silent model, dtype, attention, generation, or tokenizer fallback occurs;
\item the candidate can run the default one-process/one-active-sequence scientific worker, unless an independently justified tensor-parallel configuration has already passed its own gate.
\end{enumerate}

For each eligible candidate, execute the same frozen benchmark manifest and measure: valid logical generations per hour; generated-token throughput; end-to-end latency distribution; peak allocated and reserved device memory; technical failure rate; and monetary cost per 1,000 valid logical generations using the contemporaneous billed instance rate. The primary ranking metric is cost per 1,000 valid logical generations. Exact ties are resolved by lower technical failure rate, then higher valid generations per hour, then lexicographic frozen configuration ID. The lowest-cost eligible configuration is selected and frozen. No DCI effect estimate, verifier failure magnitude, or behavioural outcome may enter this selection.

The benchmark manifest, exact offered prices used, timing code, warm-up rule, number of benchmark generations, failure accounting, and configuration ranking script are frozen before benchmark results are inspected. Lambda bills On-Demand Cloud instances according to then-current usage pricing; current rates and billing details are volatile, so the exact price input is reverified and archived at freeze \href{https://docs.lambda.ai/public-cloud/billing/}{(Lambda, 2026d)}.

\subsection{Immutable implementation manifest}
Before any model-dependent scientific calibration, freeze at minimum: checkpoint repository/revision and every weight hash; tokenizer files and hash; exact chat-template text; OS/kernel; accelerator backend; GPU SKU; driver; CUDA or ROCm version as applicable; PyTorch and Transformers versions; attention backend and kernel settings; exact dtype; generation configuration; reasoning parser/controller; final-answer extractor; verifier commit; dependency lock; container digest; deterministic algorithm flags; device-interchangeability policy; seed-mapping code; and all behaviourally relevant environment variables. Startup fails closed on any mismatch.

PyTorch explicitly warns that exact reproducibility is not guaranteed across releases or platforms and recommends controlling randomness and deterministic algorithms within a fixed environment \href{https://docs.pytorch.org/docs/stable/notes/randomness.html}{(PyTorch, 2026a)}. Its deterministic-algorithm mode concerns repeatability for the same input on the same software and hardware and is not by itself a whole-application reproducibility guarantee \href{https://docs.pytorch.org/docs/stable/generated/torch.use_deterministic_algorithms.html}{(PyTorch, 2026b)}. NVIDIA likewise documents library reproducibility as conditional on toolkit and GPU architecture details for relevant CUDA libraries \href{https://docs.nvidia.com/cuda/cublas/index.html}{(NVIDIA, 2026)}. v19 therefore relies on empirical exact-replay certification of the complete scientific stack rather than extrapolating a library-level guarantee to the whole generator.

\subsection{Precision and quantisation}
Prefer BF16 when the selected checkpoint and every executed kernel path support BF16 correctly under the validated stack. Otherwise freeze FP16. Silent dtype fallback is prohibited. Quantisation is excluded unless a quantised checkpoint is itself explicitly defined and validated as the scientific model before the dependent calibration stages. Precision, attention backend, and relevant kernel choices are behaviourally relevant runtime fields.

\subsection{Frozen stochastic decoder}\label{sec:decoder}
Confirmatory generations use
\[
T=0.70,\qquad p_{\rm top}=0.95,\qquad k_{\rm top}=\text{disabled},
\]
with \texttt{do\_sample=true}, \texttt{num\_beams=1}, repetition penalty 1.0, \texttt{no\_repeat\_ngram\_size=0}, and no min-$p$, typical-$p$, epsilon-cutoff, eta-cutoff, frequency penalty, presence penalty, best-of-$N$, or other truncation/selection mechanism. Hugging Face exposes these controls separately and may inherit values from a generation configuration, so the complete resolved generation configuration must be frozen explicitly \href{https://huggingface.co/docs/transformers/main_classes/text_generation}{(Hugging Face, 2026a)}.

Freeze the exact representation that disables top-$k$ under the pinned Transformers version and validate that no top-$k$ logits processor is active. Greedy decoding is excluded because it defines a different decoder kernel.

\subsection{Scientific seed construction}
Each logical generation receives one scientific seed from SHA-256 over the canonical serialization of protocol ID, model hash, pair, unique superblock ID, cell condition, budget, load, presentation order, arm, and replicate. Interpret the first 64 bits as unsigned and reduce modulo $2^{63}-1$. The canonical serialization, domain separators, and mapping are frozen and collision-tested over at least 100,000 legal IDs.

The framework RNG is reset immediately before each one-sequence scientific generation. Scientific seeds are independent identifiers under the design and are never changed because an output is malformed, a worker fails, or a device is replaced.

\subsection{Hard reasoning-token controller}
A checkpoint is eligible only if its native output protocol provides an unambiguous frozen distinction between reasoning and final-answer content. The logical generation is
\[
\text{prompt}\to \text{reasoning content}\le b\to \text{final answer}\le V.
\]
Render prompts with the frozen chat template using \texttt{apply\_chat\_template(..., tokenize=True)}. Hugging Face documents that chat templates already contain model-specific control tokens and that tokenising separately with additional special tokens can duplicate them; direct templated tokenisation is therefore the reference path \href{https://huggingface.co/docs/transformers/main/chat_templating}{(Hugging Face, 2026b)}.

Count reasoning-content tokens excluding frozen control markers. If the native reasoning-end delimiter occurs before $b$, transition normally. If reasoning content reaches exactly $b$ first, stop reasoning sampling, append the frozen native reasoning-end delimiter deterministically without consuming a random draw, set \texttt{CAP\_FORCED\_TRANSITION=1}, and continue final-answer generation without resetting KV/cache state or the pseudorandom stream. Stop after exactly $V$ non-control final-answer tokens if no native terminator appears, persist that prefix, and set \texttt{FINAL\_TRUNCATED=1}. Do not sample token $V+1$.

KV/cache state may continue only across the two phases of the same logical generation. State leakage across different generation IDs is prohibited. Generic textual stop strings, best-of-$N$, reranking, self-consistency selection, critique/revision, tool-mediated replacement, and post hoc continuation are prohibited.

\subsection{Budget grid, manipulation gate, and final-answer allowance}
A broad ordered candidate budget grid $\mathcal B_{{\rm cand},m}$ is authored before manipulation outcomes are inspected. It spans a useful reasoning range while respecting the validated context envelope. On a dedicated disjoint manipulation bank, for every adjacent candidate pair $b_r<b_{r+1}$ define
\[
q_{b_r}=\frac1M\sum_{i=1}^{M}\mathbf 1\{U_{R,i}(b_{r+1})>U_{R,i}(b_r)\}.
\]
Whole prompts are resampling clusters carrying all budgets and repeats together. The prompt count $M$, repeat count, simultaneous confidence construction, and lower-bound separation threshold are numerical \PFD{} values frozen before manipulation outcomes are evaluated. $\Bret$ is the largest ordered subset whose adjacent pairs all pass the frozen criterion and with at least three budgets. Fewer than three retained budgets fails the compute-treatment gate.

$V$ is common across assigned reasoning budgets and is selected on disjoint pre-confirmatory material to provide final-answer headroom for the most demanding admissible $k=5$ compositions. A budget-dependent $V$ is prohibited. Realised reasoning use, cap hit, natural reasoning termination, final-answer length, and latency are post-treatment diagnostics rather than primary covariates.

\subsection{GPU interchangeability and failure recovery}\label{sec:interchange}
Before confirmation, every physical GPU that may execute a scientific draw belongs to one certified interchangeability class. The class requires the same frozen GPU SKU where hardware identity is relevant to replay, container image, driver/CUDA or ROCm stack, PyTorch and Transformers versions, attention backend, dtype, model and tokenizer hashes, chat template, generation configuration, controller, seed mapper, and all behaviourally relevant runtime flags.

Every device must pass the frozen cross-device replay corpus with exact agreement on token IDs, raw persisted output bytes, controller flags, final-answer extraction, and verifier outputs. A class passes only if all intended devices satisfy the same criterion. If the intended pool fails interchangeability, the full confirmatory run cannot launch with that pool. There is no fallback that permanently pins a large fraction of scientific draws to a non-interchangeable physical GPU.

When a certified device fails, the same generation may be replayed on any other currently certified device using the identical \texttt{generation\_id}, prompt bytes, scientific seed, model state, and runtime manifest. A replacement device added later must pass the complete class-certification suite before receiving any scientific draw. Retries remain attempts of one scientific observation and never create replacement observations.

Persistent storage contains the immutable scientific-draw manifest, frozen manifests, raw outputs, attempt lineage, completed-generation ledger, verifier outputs, analysis-ready records, and hashes. Ephemeral local instance storage may contain only reproducible caches and scratch data. Output bytes are durably persisted before verification is considered complete.

\subsection{vLLM is optional only}
Transformers is the initial scientific reference engine. vLLM is an optional pre-confirmatory optimisation and is never required for feasibility. Current vLLM documentation states that reproducibility is not guaranteed by default and conditions reproducible operation on explicit controls and stable hardware/version; its batch-invariance feature is also separately documented and version-dependent \href{https://docs.vllm.ai/en/latest/usage/reproducibility.html}{(vLLM Project, 2026a)}; \href{https://docs.vllm.ai/en/stable/features/batch_invariance/}{(vLLM Project, 2026b)}.

vLLM may replace Transformers only before model-dependent construction and only if it passes every Stage~2A runtime test plus a frozen paired comparison of 10,000 prompt--seed cases spanning pairs, parameter extremes, maximum context, retained/candidate budgets, both orders, all arm types, and controller boundary cases. Equivalence requires exact equality of generated token IDs, persisted output bytes, controller transition flags, extracted final answers, and deterministic verifier outputs for all 10,000 cases. One mismatch fails the gate. Failure leaves Transformers as the scientific runtime without any scientific penalty.

\subsection{Scientific interpretation}
For primitive $r$, the estimand is
\[
\tau_r=E_{X\sim P_p,\,U\sim\mathcal D_{\rm decode}}[D_r(X,U)],
\]
where $P_p$ is the frozen independent accepted-superblock distribution and $\mathcal D_{\rm decode}$ includes the checkpoint, tokenizer/template, exact stochastic sampling policy, hard reasoning/final-answer state machine, frozen runtime/backend, and deterministic seed mapping. Changing any behaviourally relevant runtime component changes the implemented intervention and invalidates dependent pre-confirmatory artifacts.

\section{Outcome measurement and verifier validation}

\subsection{Orthogonal compliance and generation-state fields}

Constraint compliance and generation state are recorded separately. The primary constraint field is
\[
\texttt{constraint\_verdict}
\in
\{\texttt{PASS},\texttt{FAIL}\}.
\]
For every delivered final-answer segment, an orthogonal generation/measurement field records one primary state from
\begin{center}
\texttt{SCORABLE}, \texttt{UNPARSEABLE}, \texttt{TRUNCATED},
\texttt{EMPTY}, or \texttt{REFUSAL}.
\end{center}
A separate Boolean field \texttt{exploit\_suspected} may coexist with any delivered state. For delivered outputs, the primary constraint verdict remains binary: an exact requirement that cannot be demonstrated as satisfied receives \texttt{FAIL}, while the state fields preserve why. A technical failure before a final-answer segment is delivered has no fabricated behavioural verdict and follows the retry/completeness rules.
Where the constraint admits a useful richer description, the verifier also records
\[
\texttt{FULL},\texttt{PARTIAL},\texttt{VIOLATION}
\]
as a secondary compliance field. For an exact target, \texttt{PARTIAL} implies primary target failure $F_A=1$.

A delivered response is not dropped because its state is malformed, refused, truncated, empty, or suspicious. $F_A=0$ only when the frozen target verifier returns PASS on the extracted final-answer segment. Technical failure before a behavioural response is accepted follows the retry policy rather than being coded as behavioural success or failure.

\subsection{Reasoning-controller instrumentation}

Every generation stores at least:
\begin{itemize}
\item assigned budget id and token cap $b$;
\item realised reasoning use $U_R$ under the frozen tokenizer;
\item whether the reasoning cap was hit;
\item whether reasoning ended naturally;
\item reasoning-parse status and final-answer-extraction status;
\item raw-reasoning-trace hash and final-answer hash;
\item checkpoint, tokenizer, chat-template, engine, dtype, and controller hashes;
\item final-answer token count, finish reason, latency, and retry lineage.
\end{itemize}
The raw reasoning trace is preserved for instrumentation and reproducibility. Primary constraint verifiers receive only the final-answer segment.

\subsection{Target, focal, and reference verifiers have different scientific roles}

A target verifier defines the primary outcome and therefore receives the strictest individual reliability requirement. A focal verifier measures focal compliance for manipulation and diagnostic analyses. Candidate/reference verifiers affect matching, satisfiability, and admission. Their errors can alter the counterfactual class even when they do not directly score $F_A$.

The audit report therefore distinguishes target-verifier reliability, focal-verifier reliability, candidate/reference-verifier reliability, pair-specific coupling, semantic wording validity, missing class support, underpowered audits, and demonstrated verifier defects. These categories are never collapsed into a single generic ``verifier passed'' flag.

\subsection{Reliability architecture}

\FIXED{} The two-tier reliability architecture from the preceding protocol is retained. The 5\% Tier-1 ceiling, 20\% Tier-2 per-stratum ceiling, pooled 5\% system certificate, and inherited 73-per-truth-class Tier-1 planning count are preserved. The exact two-sided binomial interval implementation used by the inherited gate must be recovered from the authoritative code/specification and frozen at Stage~2; if it is not uniquely recoverable, its choice is a \PFD{} within register item D9 rather than an implementer decision.

\textbf{Tier 1.} Each target and focal verifier is individually certified against a conditional error ceiling of 5\% using the frozen two-sided interval rule. The inherited planning calculation requires 73 human-adjudicated true failures and 73 human-adjudicated true passes per Tier-1 verifier unit. With three pairs and six target/focal verifier units, this alone implies 876 adjudicated labels before broader reference-verifier strata are counted.

\textbf{Tier 2.} Other implementations that can affect reference admission may use a coarser 20\% per-stratum ceiling together with a pooled 5\% system certificate. This supports a coarser system-level claim and does not imply that every reference verifier individually satisfies the Tier-1 5\% ceiling.

A missing denominator is missing evidence, not zero error. A false-pass bound requires adjudicated examples that truly fail the constraint. A false-fail bound requires adjudicated examples that truly pass. If either required class is absent, the corresponding reliability bound is unavailable and the unit cannot pass that gate.

\subsection{Human and automated audit requirements}

Human adjudication is blinded to arm, automatic verdict, and constraint role where feasible. Audit material includes boundary cases, parameter extremes, realistic malformed outputs, chat preambles, truncation, empty output, Unicode normalisation, zero-width characters, homoglyphs, non-breaking spaces, smart quotes, combining marks, code fences, and line-ending variants.

Automated suites include differential tests, property-based tests, metamorphic tests, and pair-specific structural tests. A signed audit cannot be overwritten by a later run. A demonstrated semantic or verifier defect requires repair and re-audit before freeze; collecting more confirmatory data is not a remedy for invalid measurement.

\subsection{Wording policy and semantic audit}

Confirmation uses one canonical wording per constraint identity. Held-out renderings are reserved for wording robustness and reference-comparability validation. Every rendering must preserve quantifier, threshold, units, scope, relational direction, and exact/minimum/maximum semantics.

Canonical wordings are part of the instruction-level treatment and are frozen with hashes. The same canonical wording for a given constraint identity is used across paired loads, budgets, orders, and applicable arms. A wording that changes the semantic rule is a different intervention and cannot be treated as a paraphrase.



\subsection{Behavioural outcomes versus technical failures}
A generation is scientifically observed when model output bytes are successfully produced and persisted. Refusal, malformed output, an empty delivered answer, instruction failure, and truncation caused by the frozen $b$ or $V$ caps are behavioural outcomes. They receive ordinary behavioural/verifier states under the frozen rules.

Technical missingness is narrower: the intended scientific output cannot be recovered because computation failed before a valid output record was persisted, for example unrecoverable worker death, hardware or driver failure, OOM before persistence, or unrecoverable transport corruption. If generation succeeded and only deterministic verification failed operationally, rerun the verifier on the stored bytes; never regenerate the model output.

\subsection{Implementation-readiness snapshot}
Implementation status is a pre-confirmatory engineering fact, not scientific evidence for DCI. At the v19 drafting date, the provenance, retry-lineage, frozen-manifest, causal-field-role, completeness, and substantial reference-construction machinery are reusable subject to v19 hash/semantic validation. Pair~2 mandatory-word and all-words-unique verifiers exist. Pair~3 exact-sentence-count and increasing-length verifiers exist. Pair~1 bullet-count and transitive-order verifier support remains to be built. The bounded/partial-identification inference layer, hard reasoning-token controller, and superblock identity/order/paired-load spine require implementation. These statements are recorded only to identify gates that must be passed before confirmation; passing software tests does not constitute evidence for the DCI hypotheses.

\section{Tasks, superblock identity, randomisation, and provenance}

\subsection{Disjoint material splits}

All material that can affect a frozen scientific choice is partitioned before confirmation. At minimum the protocol uses the following mutually disjoint task/topic pools, with disjointness enforced at the most conservative content-identity level available:
\begin{enumerate}
\item \textbf{development}: software development, verifier unit tests, prompt debugging, and exploratory adapter work;
\item \textbf{parameter calibration}: estimation of admissible parameter ranges $\mathcal P_Q$ and sampling distributions $\nu_Q$;
\item \textbf{reference construction}: calibration of focal and candidate matching signatures, sensitivity coefficients, calipers, and candidate admission quantities;
\item \textbf{independent reference validation}: held-out confirmation that admitted references satisfy the frozen comparability and semantic requirements;
\item \textbf{confirmation}: the superblocks that enter the frozen DCI estimands;
\item \textbf{terminal falsification}: a sequestered final slice used only after the confirmatory analysis has been locked.
\end{enumerate}
A separate adapter/manipulation prompt bank may be used for hard-budget validation when it is demonstrably disjoint from all DCI content pools. No confirmatory task, wording, answer, or model output may enter parameter calibration, reference construction, adapter selection, power-scenario tuning, or other pre-freeze design decisions.

\subsection{Content-sufficiency gate}

Before expensive calibration begins, the study must establish that every target--focal pair has enough disjoint material for all six scientific pools above and for every generative scope required by that pair. The sufficiency report must count, separately by pair and scope:
\begin{itemize}
\item eligible tasks/topics;
\item focal and target wordings;
\item candidate reference identities;
\item unrelated targets needed for incoming-vulnerability estimation;
\item parameter-calibration instances;
\item human-audit examples with both true pass and true fail support;
\item feasible background candidates;
\item independent validation and terminal-falsification material.
\end{itemize}
A pair that cannot populate all required pools without reuse is a pre-freeze feasibility failure. Material may be authored before freeze under a versioned content procedure; once a pool is signed and hashed, later movement of items between pools is prohibited.

\subsection{Identity hierarchy}
The implementation distinguishes stochastic scientific draws from execution attempts. For pair $p$, accepted draw index $i$ and its pretreatment content define a \texttt{context\_fingerprint}. The unique statistical-unit identifier is
\[
\texttt{superblock\_id}=\operatorname{hash}(\texttt{protocol},p,i,\texttt{context\_fingerprint}).
\]
Including $i$ is mandatory because with-replacement sampling may produce identical contexts at different independent draw indices.

A cell is
\[
\texttt{cell\_id}=\operatorname{hash}(\texttt{superblock\_id},b,k,o,q),
\]
where $q\in\{B,R_1,R_2,R_3,R_4\}$. A planned stochastic replicate is
\[
\texttt{generation\_id}=\operatorname{hash}(\texttt{cell\_id},g),\qquad g\in\{1,\ldots,G\}.
\]
A technical execution attempt is
\[
\texttt{attempt\_id}=\operatorname{hash}(\texttt{generation\_id},a),
\]
where $a$ is the retry index. Attempts never create new scientific observations.

\subsection{Presentation order as a crossed primary moderator}
\FIXED{} every superblock is rendered in
\[
o_A:\ A\text{ before added-slot occupant},\qquad
o_Q:\ \text{added-slot occupant before }A.
\]
The three background constraints retain their frozen relative order and prompt positions. Changing order swaps only the target and added-slot positions under the frozen template. The primitive is $\tau_{p,j}(b,k,o)$, not an order average. Direct order contrasts are secondary-confirmatory as specified in Section~\ref{sec:orderfamily}.

\subsection{Replicates}
For each causal cell, $G$ is the number of independent stochastic generation replicates under the frozen decoder. $G$ is selected prospectively with pair-specific $H_p$. Replicates differ only by their replicate identifier and corresponding seed. They estimate decoder expectation within a superblock and are averaged before claim-bearing inference. Increasing $G$ does not increase the number of independent superblocks.

\subsection{Randomisation, scheduling, and temporal interleaving}
A versioned root seed and domain-separated seed mapping implement all construction and scheduling randomness. Superblock proposal streams are separated by pair and accepted draw index as Section~\ref{sec:bgsampling} requires. The execution scheduler interleaves pair, superblock, budget, load, order, arm, and replicate so that calendar time and machine state are not systematically aligned with treatment. Scheduling order cannot alter scientific IDs or seeds.

Every generation starts from a fresh model context. No KV cache, conversation state, tool state, hidden scratch state, or prior sampled text may leak across generation IDs. The default one-active-sequence-per-GPU architecture makes scientific draws operationally independent. Any later batching optimisation must pass the exact equivalence gate before it can become scientific runtime.

\subsection{Causal field-role registry}
Every analysis field has exactly one pre-frozen role:
\begin{itemize}
\item \textbf{identity}: protocol, pair, accepted draw index, context fingerprint, superblock, cell, generation, and attempt IDs;
\item \textbf{treatment}: added-slot arm, assigned hard budget $b$, load $k$, and order $o$;
\item \textbf{pretreatment design}: task, target, focal/reference identity, parameters, wordings, $S^{\max}$, prompt template, sampling distributions, and construction fields;
\item \textbf{apparatus/provenance}: checkpoint revision, tokenizer/template, controller, verifier, inference, container, accelerator/backend, GPU class, driver/runtime, dtype, attention backend, and manifest hashes;
\item \textbf{post-treatment instrumentation}: realised reasoning use $U_R$, cap hit, natural termination, final-answer length, latency, finish reason, and truncation state;
\item \textbf{outcome}: target verdict and $F_A$;
\item \textbf{diagnostic}: focal/reference compliance, richer partial-compliance state, exploit flags, and prespecified non-primary summaries.
\end{itemize}
Primary estimators reject post-treatment and diagnostic fields as covariates, weights, filters, matching coordinates, or denominator definitions. Realised reasoning use and answer length may be described or studied secondarily but cannot redefine the causal replacement contrast \href{https://doi.org/10.1146/annurev-publhealth-032315-021402}{(VanderWeele, 2016)}.

\subsection{Immutable provenance and identical-draw retries}
Every planned generation has immutable scientific identity and separate attempt lineage. A retry may reuse only the identical prompt bytes, checkpoint/model state, tokenizer/template, decoder, hard budget, $V$, scientific seed, controller, runtime manifest, and \texttt{generation\_id}. A new seed, edited prompt, alternate continuation, or replacement generation is prohibited.

If computation fails before output bytes are durably persisted, replay may occur on any currently certified member of the frozen interchangeability class in Section~\ref{sec:interchange}. The failed physical device has no permanent ownership of the draw. If no certified device remains available, retries pause until a replacement device passes the same class gate. After the frozen retry policy is exhausted, the scientific generation becomes permanently technically missing and enters interval completion. Behavioural outputs that are malformed, refused, truncated, empty, or scientifically inconvenient are never retried merely because of their content.

The persistent ledger stores, for every generation, planned identity, every attempt, device identity, start/end time, failure category, output hash/path when successful, and terminal observed/missing state. Deterministic verification can be replayed from stored output bytes without regenerating the model.

\section{Primary estimands and scientific hypotheses}
\label{sec:estimands}

\subsection{Superblock-level primitive contrast}

For pair $p$, let $i\in\{1,\ldots,H_p\}$ index its frozen independent confirmatory superblocks. Let $q$ identify either the focal arm $B$ or one admitted reference $R_j$. Define the replicate mean target failure
\[
\bar Y_{i,p,q}(b,k,o)
=
\frac{1}{G}\sum_{g=1}^{G}F_{A,i,p,q,b,k,o,g}.
\]
The superblock-level reference-specific DCI contrast is
\[
D_{i,p,j}(b,k,o)
=
\bar Y_{i,p,B}(b,k,o)
-
\bar Y_{i,p,R_j}(b,k,o).
\]
With complete observations, the descriptive estimator is the unweighted mean across that pair's independent superblocks:
\[
\widehat\tau_{p,j}(b,k,o)
=
\frac{1}{H_p}\sum_{i=1}^{H_p}D_{i,p,j}(b,k,o).
\]
If permanent technical missingness occurs, the point-estimation convention in Section~\ref{sec:analysis} applies while claim-bearing inference uses interval completion. Pair-specific $H_p$ changes precision only. The focal/reference pairing is preserved, and no primary reference weights $w_j$ exist.

\subsection{Population interpretation}
For pair $p$, $\tau_{p,j}(b,k,o)$ averages the focal-minus-reference superblock contribution over the frozen accepted-superblock distribution and frozen stochastic decoder. The pair-specific proposal law draws tasks with replacement from $\mathcal T_p$, one background identity with replacement from each frozen stratum at $k=5$, family-specific parameters from frozen $\nu_Q$, and the frozen position permutations; it then conditions on the pre-confirmatory whole-composition acceptance event. This population is therefore defined by scientific construction rules before confirmatory outcomes exist.

For fixed primitive $r$ within pair $p$, accepted superblocks satisfy
\[
D_{1r},\ldots,D_{H_pr}\ \text{i.i.d.},\qquad E[D_{ir}]=\tau_r,
\]
under the prespecified inferential model. Equivalently, the betting condition $E[D_{ir}\mid\mathcal F_{i-1}]=\tau_r$ holds. Pair identity remains fixed and descriptive rather than an exchangeable random population factor. Unequal $H_p$ affects precision, not the estimand.

\subsection{Reference robustness}

For each retained cell $s=(p,b,k,o)$,
\[
R_s=R_p(b,k,o)=\min_{j\in\mathcal J_p}\tau_{p,j}(b,k,o).
\]
A positive $R_s$ means the focal constraint raises target failure relative to every one of the four admitted references. If the primitive effects differ in sign, the scientific result is reference-dependent DCI rather than a single averaged effect.

\subsection{Critical-interference headline}

The primary global headline is
\[
\Tcrit
=
\max_{s\in\Sc}R_s
=
\max_{(p,b,k,o)\in\Sc}
\min_{j\in\mathcal J_p}\tau_{p,j}(b,k,o).
\]
The search set $\Sc$ is frozen after all pre-freeze support and manipulation gates and before confirmatory generation. No cell may enter or leave $\Sc$ because its confirmatory effect is large, small, significant, or surprising.

The two operators have different scientific meanings. The inner $\min_j$ asks whether the effect survives the least favourable admissible counterfactual. The outer $\max_s$ asks whether any preregistered region exhibits a sharp reference-robust interference effect. Because the global headline searches across many cells, the primitive bounds must be simultaneous over the entire searched primitive family.

\subsection{Compute dependence is reference-specific}

Let $b_{\min}$ and $b_{\max}$ denote the smallest and largest retained hard reasoning caps for the selected checkpoint. For each pair, reference, load, and order,
\[
\Delta_b\tau_{p,j}(k,o)
=
\tau_{p,j}(b_{\max},k,o)
-
\tau_{p,j}(b_{\min},k,o).
\]
Positive values mean that reference-specific interference is larger at the highest retained reasoning cap; negative values mean it is smaller. Budget is categorical. The design imposes no monotonicity assumption and no linear per-token interpretation. Adjacent-budget contrasts are secondary.

Raw token counts are meaningful only under a fixed tokenizer and model adapter. Compute contrasts are interpreted within a checkpoint. A 2,048-token cap under one tokenizer is not treated as the same cardinal compute dose as a 2,048-token cap under another model family.

\subsection{Load dependence is reference-specific and paired}

Because both loads arise from the same superblock,
\[
\Delta_k\tau_{p,j}(b,o)
=
\tau_{p,j}(b,5,o)
-
\tau_{p,j}(b,2,o)
\]
is a paired load contrast. Positive values mean that the focal-versus-reference interference contrast is larger when the same scientific context is augmented with the feasible three-constraint background set $S^{\max}$.

No independent $k=2$ and $k=5$ block populations are constructed. Every paired load contribution is formed within the same superblock. If the non-claim-bearing diagnostic bootstrap is run, it resamples the whole superblock so both loads remain paired.

\subsection{Pair identity is descriptive, not a random factor}

The three target--focal pairs are deliberately selected exemplars. They are not exchangeable random draws from a population of constraint pairs. Pair identity is perfectly confounded with the chosen target-family and focal-family combination. Consequently, no scalar pair-heterogeneity hypothesis is primary. Cross-pair differences in $\{\tau_{p,j}(b,k,o)\}$ are reported descriptively.

Population-level claims about constraint families require a future design with multiple independently sampled target--focal pairs per family combination.

\subsection{SESOI}\label{sec:sesoi}
Let $\epsilon$ denote the absolute smallest effect size of scientific interest on the target-failure probability scale. \FIXED{} at Stage~1,
\[
\epsilon\ge0.06.
\]
This lower bound is itself a scientific commitment. It ensures that the already fixed construction ceilings resolve numerically as
\[
\min(0.015,\epsilon/4)=0.015,
\qquad
\min(0.02,\epsilon/2)=0.02.
\]
The exact value of $\epsilon$ remains a \PFD{} only until the earliest procedure that numerically depends on that exact value. In particular, if the 2,000-observation verifier-coupling bank fails to certify and the $\epsilon/8$ precision expansion branch is invoked, exact $\epsilon$ must be frozen before any expansion observations are generated. Exact $\epsilon$ must also be frozen before reference-matching constants or any other construction quantity is calculated from it, before inference-method comparison uses scenarios at $\Theta=\epsilon$ or $2\epsilon$, and before authoritative $(\mathbf H,G)$ selection. Once exact $\epsilon$ has been used by any such procedure it cannot be changed without invalidating all dependent artifacts.

\paragraph{Prospective exact-SESOI decision procedure.}
Before Stage~9, and earlier whenever an exact-$\epsilon$ dependency is reached, the investigators must complete the following scientific-decision procedure before any exact value is consumed:
\begin{enumerate}
\item write and sign a short SESOI rationale before executing the dependent procedure;
\item state why the selected value, subject to the fixed floor $\epsilon\ge0.06$, is the smallest absolute increase in target-failure probability that would be scientifically meaningful for this pilot;
\item record every substantive precedent, prior empirical result, domain judgement, or practical interpretation used to choose the value;
\item explain why effects smaller than the selected value would not support the SESOI-level scientific claim;
\item select the exact value without consulting inference-method comparison results, prospective power curves generated at alternative candidate values of $\epsilon$, Stage~10 attainable-bank results, or confirmatory outcomes; and
\item enter the exact value, signed rationale, source record, and freeze timestamp into the signed pre-freeze decision record before any procedure numerically consuming exact $\epsilon$ is executed.
\end{enumerate}
This procedure documents a prospective scientific judgement; it does not estimate $\epsilon$ from design-optimisation outputs. An earlier exact dependency, including the $\epsilon/8$ verifier-coupling expansion, moves this procedure earlier in the stage order. The already fixed construction ceilings above remain $0.015$ and $0.02$ for every admissible exact $\epsilon$ because the Stage~1 floor is unchanged.

The SESOI is a scientific importance threshold. A minimum detectable effect is a prospective property of a design and is never substituted for $\epsilon$ after outcomes are observed \href{https://doi.org/10.1177/1948550617697177}{(Lakens, 2017)}.

\section{Statistical analysis, simultaneous inference, and claim logic}\label{sec:analysis}

\subsection{Analysis sequence}
The confirmatory analysis is locked before the first confirmatory generation. It: (1) verifies the immutable manifest, runtime, family cardinalities, and identity invariants; (2) assigns structural failure states without consulting effect direction; (3) maps every planned generation to an observed binary target-verifier value or a support interval after exhausted technical retries; (4) aggregates generation replicates into bounded superblock-level primitive intervals; (5) applies the single inference procedure selected and frozen at the pre-confirmatory inference-method qualification stage; (6) constructs simultaneous bounds within each prespecified inferential family; (7) derives least-favourable-reference and global max--min bounds deterministically; and (8) emits the frozen result taxonomy. No confirmatory bootstrap, model switch, hardware switch, sample-size extension, or method re-selection occurs after outcomes are available.

\subsection{Global primary primitive family}
For every retained pair, reference, hard budget, load, and order, define one primitive index
\[
r=(p,j,b,k,o).
\]
With three pairs, four references, two loads, two orders, and $|\Bret|$ retained budgets, the complete primary one-sided family has
\[
m_{\rm prim}=3\times4\times2\times2\times|\Bret|=48|\Bret|.
\]
This family is protected simultaneously at $\alpha_{\rm prim}=0.05$. No primitive is removed because its observed variance is zero, its point estimate is small, its interval is wide, or its cell is not ultimately selected by the outer maximum.

\subsection{Exact bounded supports}
For observed binary target failures, the replicate mean in any arm lies in $[0,1]$. Therefore an ordinary focal-minus-reference superblock contribution satisfies
\[
D_{ir}\in[-1,1],\qquad a_r=-1,\quad b_r=1,\quad w_r=b_r-a_r=2.
\]
This support is exact and is used even when all realised observations occupy a narrower range.

Every secondary-confirmatory compute, load, or order contribution is the difference of two primitive contributions. Therefore
\[
Q_{ir}\in[-2,2],\qquad a_Q=-2,\quad b_Q=2,\quad w_Q=4.
\]
The support width $w_Q=4$ is not inferred from observed data. The exact interval-completion endpoints for each family are given in Section~\ref{sec:secondarysupports}.

\subsection{Inference candidates and pre-confirmatory qualification}
Exactly two confirmatory confidence procedures are candidates in v19:
\begin{enumerate}
\item \textbf{MP-EB}: the Maurer--Pontil empirical-Bernstein construction retained as the conservative baseline, extended conservatively to interval-valued contributions as specified below;
\item \textbf{WSR-PP}: a one-sided positive-capital predictable-plug-in betting construction based on the bounded-mean martingale framework of Waudby-Smith and Ramdas, with upper bounds obtained by reflection and two-sided intervals obtained by simultaneous lower/upper protection.
\end{enumerate}
No third method may be introduced after the comparison data are inspected. A new method requires a new protocol version or an explicitly frozen pre-comparison amendment made before any method-comparison output exists.

The inference-method qualification stage is \textbf{validity first}. Each candidate must pass: (i) a written theorem/assumption mapping to the exact DCI superblock law; (ii) exact-support tests; (iii) one-sided and two-sided coverage derivations; (iv) zero-variance fixtures; (v) interval-missingness proofs and brute-force small-sample checks; (vi) complete-family multiplicity tests; (vii) max--min transformation tests; (viii) null-coverage simulation; and (ix) agreement between production code and an independently implemented reference calculator on a frozen fixture bank. A method failing any validity gate is ineligible regardless of interval width or compute savings.

Among valid candidates, the frozen 10,000-experiment comparison in Section~\ref{sec:power} estimates prospective power, interval width, required $H_p$, and expected compute under the same scenarios and family sizes. The deterministic selection rule chooses the valid method with the lowest expected-compute design satisfying all method-comparison acceptance criteria. If exact expected compute ties, choose the design with smaller $\sum_pH_p$; if still tied, choose the method with smaller median primary interval width under the central $\Theta=\epsilon$ scenario; if still tied, choose lexicographically by frozen method ID. If WSR-PP fails any validity gate, MP-EB remains primary. If no valid betting method exists, no claim that betting improves feasibility is made. If MP-EB is valid but no method yields an adequately powered feasible design, MP-EB is frozen as the conservative primary method and the study follows the feasibility-and-measurement-pilot fallback. If neither candidate passes statistical validity, confirmation cannot begin.

After selection, freeze the method ID, formulas, all numerical tuning parameters, primitive ordering, numerical inversion tolerance, multiplicity allocation, missingness implementation, software environment, production code hash, independent reference implementation hash, and validation report. Method selection occurs before authoritative attainable-bank $(\mathbf H,G)$ selection and before any confirmatory outcome.

\subsection{Candidate 1: Maurer--Pontil empirical-Bernstein}
For complete primitive contributions $D_{1r},\ldots,D_{H_rr}\in[a_r,b_r]$, define
\[
\bar D_r=\frac1{H_r}\sum_{i=1}^{H_r}D_{ir},
\qquad
V_r=\frac{1}{H_r(H_r-1)}\sum_{1\le i<j\le H_r}(D_{ir}-D_{jr})^2.
\]
For a one-sided primary family of size $m_{\rm prim}$, set
\[
\lambda_{\rm MP}=\log\!\left(\frac{2m_{\rm prim}}{\alpha_{\rm prim}}\right).
\]
The candidate lower bound is
\[
L_r^{\rm MP}
=\max\!\left\{a_r,
\bar D_r-\sqrt{\frac{2V_r\lambda_{\rm MP}}{H_r}}
-\frac{7w_r\lambda_{\rm MP}}{3(H_r-1)}\right\},
\qquad H_r\ge2.
\]
This is the finite-class union-bound specialization of the bounded empirical-Bernstein result of Maurer and Pontil \href{https://arxiv.org/abs/0907.3740}{(Maurer \& Pontil, 2009)}. For primary primitives $w_r=2$. $V_r=0$ is an ordinary numerical value; it removes the variance term but leaves the finite-sample range term. There is no standard-error denominator and no zero-variance failure state.

For a two-sided family of size $M$, use
\[
\lambda_{2,\rm MP}=\log\!\left(\frac{4M}{0.05}\right)
\]
in the same lower-bound formula and its sign-reflected upper-bound analogue. This corresponds to simultaneous protection of both tails for all $M$ family members by the fixed union allocation.

The MP-EB range term depends on $H_r$ and $w_r$ but not on decoder replicate count $G$. Consequently, increasing $G$ cannot substitute for independent superblocks when that term is the power bottleneck. Stage~11 authoritative design simulations must represent this fact directly rather than treating $H$ and $G$ as interchangeable compute units.

\subsection{Candidate 2: WSR-PP one-sided betting lower bound}\label{sec:wsr}
The scientific reason for the betting candidate is to obtain a finite-sample bounded-mean procedure whose validity does not require studentisation and whose interval width may be materially smaller than MP-EB at moderate $H$. Waudby-Smith and Ramdas show that for $[0,1]$-valued variables satisfying $E[X_t\mid\mathcal F_{t-1}]=\mu$, capital processes of the form $\prod_t[1+\lambda_t(\mu)(X_t-\mu)]$ are nonnegative test martingales for allowable predictable stakes \href{https://doi.org/10.1093/jrsssb/qkad009}{(Waudby-Smith \& Ramdas, 2024)}. v19 uses only this ordinary bounded-mean result because Section~\ref{sec:bgsampling} prospectively makes the superblock process conditionally mean-stable; it does not assume the paper's separate finite-population without-replacement theorem applies to DCI.

For any contribution $D_{ir}\in[a_r,b_r]$ with width $w_r>0$, normalize
\[
X_{ir}=\frac{D_{ir}-a_r}{w_r}\in[0,1].
\]
Fix the superblock order by ascending accepted draw index $i=1,\ldots,H_r$ before any confirmatory outcome is viewed. For a family of size $M$ with desired one-sided familywise error $\alpha_F$, allocate
\[
\delta=\frac{\alpha_F}{M}
\]
to each member. This is the generic WSR-PP family-level rule. For the primary one-sided primitive family specifically,
\[
\alpha_F=\alpha_{\rm prim}=0.05,
\qquad
M=M_{\rm prim}\equiv m_{\rm prim}=48|\Bret|,
\]
and therefore
\[
\boxed{\delta_{\rm prim}=\frac{0.05}{48|\Bret|}}.
\]
Other WSR-PP inferential families retain the same generic rule with their own prespecified familywise level, cardinality, and, for two-sided families, tail allocation. For testing a candidate mean $m\in[0,1]$, construct the smoothed predictable summaries from prior lower-endpoint observations $x_1^-,\ldots,x_{t-1}^-$:
\[
\widehat\mu_t=\frac{1/2+\sum_{i=1}^{t}x_i^-}{t+1},
\qquad
\widehat\sigma_t^2=\frac{1/4+\sum_{i=1}^{t}(x_i^- - \widehat\mu_i)^2}{t+1},
\]
with $\widehat\sigma_0^2=1/4$. Define the fixed-horizon predictable base stake
\[
\widetilde\lambda_t
=\sqrt{\frac{2\log(2/\delta)}{H_r\widehat\sigma_{t-1}^2}},
\]
and use truncation constant $c=1/2$:
\[
\lambda_t^+(m)=\min\!\left\{\widetilde\lambda_t,\frac{c}{m}\right\},
\]
where $c/m=+\infty$ when $m=0$. The truncation branch is operationally material rather than a numerical edge case. For illustration, at the minimum retained budget count $|\Bret|=3$, the primary allocation is $\delta=0.05/144$, and at $H_r=200$ with $\widehat\sigma_{t-1}^2=0.0225$,
\[
\widetilde\lambda_t
=
\sqrt{\frac{2\log(5760)}{200(0.0225)}}
\approx1.962.
\]
At candidate mean $m=0.55$, $c/m=0.5/0.55\approx0.909$, so the truncation bound determines the operative stake. Larger $|\Bret|$ weakly increases the base stake through the smaller memberwise $\delta$. These values are illustrative validation targets only and cannot be used to retune $c$. The one-sided positive capital based on the available lower endpoints is
\[
K_{H_r}^+(m)=\prod_{t=1}^{H_r}\left[1+\lambda_t^+(m)(x_t^- - m)\right].
\]
All factors are nonnegative because $0\le x_t^-\le1$ and $\lambda_t^+(m)\le c/m<1/m$ for $m>0$. The base stake depends only on earlier observations and the frozen horizon, so it is predictable. The smoothing constants, truncation $c$, logarithmic factor, update order, floating-point precision, log-domain accumulation algorithm, and root-finding tolerance are part of the frozen method specification.

For complete data, $x_t^-=X_{tr}$. The normalized lower confidence bound is the upper edge of the rejected low-mean set,
\[
L_{X,r}^{\rm WSR}
=\sup\left(\{m\in[0,1]:K_{H_r}^+(m)\ge1/\delta\}\cup\{0\}\right).
\]
Because every factor is nonincreasing in $m$ under the specified truncation, the rejection set is a lower interval. Compute the boundary in log space by deterministic bisection on $[0,1]$ until the enclosing interval width is at most $2^{-40}$; return the lower bisection endpoint so numerical error cannot increase the reported lower bound. Transform back as
\[
L_r^{\rm WSR}=a_r+w_rL_{X,r}^{\rm WSR}.
\]
At the true normalized mean, the latent complete-data capital is a nonnegative test martingale. Markov/Ville control at the fixed horizon gives memberwise error at most $\delta$, and the union bound gives familywise error at most $\alpha_F$. The exact code and theorem mapping are nevertheless required to pass the pre-confirmatory validation gate before WSR-PP is eligible.

\subsection{Permanent technical missingness by interval completion}\label{sec:missingness}
Technical missingness is handled without a parametric MCAR/MAR model and without assuming that missingness for the current superblock is independent of that superblock's latent outcome or generated content, conditional on the prospective complete-data superblock/filtration property stated below. For every superblock contribution, first derive its exact feasible interval $[\ell_{ir},u_{ir}]\subseteq[a_r,b_r]$ by propagating each permanently missing binary generation as $[0,1]$ through the frozen replicate average and contrast algebra. Observed behavioural failures remain observed $0/1$ values; only unrecoverable technical absence creates an interval.

For MP-EB, use the lower mean
\[
\bar\ell_r=\frac1{H_r}\sum_i\ell_{ir}
\]
and conservative pairwise variance envelope
\[
V_r^+=\frac{1}{H_r(H_r-1)}\sum_{i<j}
\max\{(\ell_{ir}-u_{jr})^2,(u_{ir}-\ell_{jr})^2\}.
\]
Then
\[
L_{r,\rm miss}^{\rm MP}
=\max\!\left\{a_r,
\bar\ell_r-\sqrt{\frac{2V_r^+\lambda_{\rm MP}}{H_r}}
-\frac{7w_r\lambda_{\rm MP}}{3(H_r-1)}\right\}.
\]
Pointwise, $\ell_{ir}\le D_{ir}$ and $V_r^+$ upper-bounds the latent pairwise sample variance, so this is no larger than the corresponding complete-data MP lower bound.

For WSR-PP, normalize the interval endpoints
\[
x_{ir}^-=\frac{\ell_{ir}-a_r}{w_r},\qquad
x_{ir}^+=\frac{u_{ir}-a_r}{w_r}.
\]
Use $x_{ir}^-$ in the positive-capital construction, including the predictable stake updates. For the same realised predictable stake sequence and any candidate $m$, $x_{ir}^-\le X_{ir}$ implies
\[
1+\lambda_i^+(m)(x_{ir}^- -m)\le
1+\lambda_i^+(m)(X_{ir}-m).
\]
The stake sequence computed from prior lower endpoints is measurable with respect to the latent past filtration because those observed lower endpoints are functions of past scientific records. Therefore the latent process using those same predictable stakes remains a valid test martingale at the true mean, while the interval-completed capital is pathwise no larger. A false lower rejection can therefore occur no more often than under the latent complete data.

The required prospective condition is on the latent complete-data sequence. Let $X_t$ be the normalized complete-data contribution that would be observed for superblock $t$ absent technical missingness, and let $\mathcal F_{t-1}$ contain the complete past scientific record available to the predictable stake rule, including prior technical-missingness states and interval endpoints. The frozen superblock law must continue to satisfy
\[
E[X_t\mid\mathcal F_{t-1}]=\mu.
\]
Current-superblock missingness may depend arbitrarily on that superblock's latent outcome or generated content because interval completion is pathwise conservative. Past missingness states, however, must not carry information that changes the conditional distribution or conditional mean of future latent superblock outcomes in a way that violates the frozen superblock law. The preregistered independent/i.i.d. superblock construction, per-superblock hardware-failure handling, frozen runtime, and identical-draw retry architecture are intended to supply this property. Thus no parametric MCAR/MAR model and no assumption of contemporaneous outcome-independent missingness is required; validity remains conditional on the prospective superblock/filtration property above.

For an upper confidence bound under WSR-PP, apply the same lower-bound algorithm to reflected latent variables $Z_i=1-X_i$. The support interval for $Z_i$ is $[1-x_i^+,1-x_i^-]$, so its conservative lower endpoint is $z_i^-=1-x_i^+$. If $L_Z$ is the resulting lower bound for $E[Z]$, set
\[
U_X=1-L_Z,\qquad U_D=a_r+w_rU_X.
\]
This reflection rule is used for every WSR-PP two-sided family. Zero empirical variance, all-zero observations, all-one observations, and all-interval observations have explicit fixtures and do not trigger a special inferential failure branch.

\subsection{Point estimates under interval-valued missingness}
For descriptive reporting, each primitive with interval-valued superblock contributions reports the lower-endpoint mean $\bar\ell_r$, upper-endpoint mean $\bar u_r$, midpoint mean, and complete-case mean where defined. No one of these descriptive conventions replaces the claim-bearing confidence bound. The interval span is displayed whenever permanent technical missingness occurs.

\subsection{From simultaneous primitive bounds to the global headline}
Let $L_r$ denote the lower bound from the single selected primary method. For cell $s$, define
\[
L_s=\min_{j\in\mathcal J_p}L_{p,j}(b,k,o),
\qquad
L_\Theta=\max_{s\in\Sc}L_s.
\]
On the event that every primary primitive lower bound covers, $L_r\le\tau_r$ simultaneously. Hence $\min_jL_{j,s}\le\min_j\tau_{j,s}$ for every cell, and taking the maximum over the same frozen $\Sc$ yields
\[
L_\Theta\le\Tcrit.
\]
No additional multiplicity correction is needed for these deterministic transformations because simultaneous primitive coverage has already protected the entire searched family.

\subsection{Selected descriptive maximum and tie-breaking}
The plug-in quantity
\[
\widehat\Theta_{\rm crit}=\max_{s\in\Sc}\min_j\widehat\tau_{j,s}
\]
is a selected descriptive estimate. Its maximising cell $s^\ast$ is selected from noisy estimates and is susceptible to winner's-curse bias. It is not the claim-bearing lower bound. Exact ties use a frozen lexicographic ordering over pair ID, budget order, load, and presentation order. Every report places $\widehat\Theta_{\rm crit}$, $s^\ast$, the selected primitive estimates, and $L_\Theta$ adjacent to the warning that the selected plug-in magnitude can be upward biased \href{https://doi.org/10.1093/qje/qjad043}{(Andrews, Kitagawa, \& McCloskey, 2024)}.

\subsection{Primary global claims}
The primary family controls its own familywise error rate at 0.05. The strongest positive headline is supported only if $L_\Theta>0$. A scientifically important positive headline at the SESOI is supported only if $L_\Theta>\epsilon$. If $L_\Theta\le0$, the study does not conclude absence of DCI; it concludes that the prespecified simultaneous lower-bound procedure did not establish a positive reference-robust regime. If the prospective design was frozen under the feasibility-and-measurement-pilot fallback, positive-looking descriptive estimates cannot be promoted to the powered headline claim.

\subsection{Multiplicity policy across confirmatory families}\label{sec:multiplicity}
\FIXED{} v19 makes \emph{separate-family} error guarantees. The primary one-sided primitive family controls FWER at $\alpha=0.05$. Each secondary-confirmatory family below separately controls its own FWER at $0.05$ under its own cardinality and two-sided allocation. These guarantees do not imply, and the report must not state, a single 5\% familywise error rate jointly across every primary and secondary claim in the experiment.

The two-sided primitive family, highest-minus-lowest compute family, paired load family, and direct paired order family are distinct confirmatory families. The order family is additionally hierarchically activated: confirmatory order-effect claims are interpreted only if the primary positive-headline gate $L_\Theta>0$ is met. This hierarchy restricts interpretation but is not used to claim a stronger across-all-families 5\% guarantee.

\subsection{Upper and two-sided primitive bounds}
A separate two-sided primitive family contains the same $M_{\rm 2prim}=48|\Bret|$ primitive means. Under MP-EB, use $\lambda_{2,\rm MP}=\log(4M_{\rm 2prim}/0.05)$ for each tail. Under WSR-PP, allocate tail probability
\[
\delta_{\rm tail}=\frac{0.05}{2M_{\rm 2prim}}
\]
and obtain lower and upper bounds using the positive-capital and reflection constructions. Two-sided primitive intervals are used for non-headline directional reporting and do not change primary lower bounds.

\subsection{Compute, load, and order families with exact endpoints}\label{sec:secondarysupports}
Let a primitive superblock interval at condition $c$ be $[\ell_i(c),u_i(c)]\subseteq[-1,1]$.

\textbf{Highest-minus-lowest compute.} For each pair, reference, load, and order,
\[
Q_i^{(b)}=D_i(b_{\max})-D_i(b_{\min})\in[-2,2],
\]
with interval completion
\[
[\ell_i^{(b)},u_i^{(b)}]
=
[\ell_i(b_{\max})-u_i(b_{\min}),\;u_i(b_{\max})-\ell_i(b_{\min})].
\]
There are
\[
M_b=3\times4\times2\times2=48
\]
confirmatory compute contrasts. Adjacent-budget contrasts are prespecified secondary descriptive intervals unless separately included in a frozen secondary family before confirmation.

\textbf{Paired load.} For each pair, reference, retained budget, and order,
\[
Q_i^{(k)}=D_i(k=5)-D_i(k=2)\in[-2,2],
\]
with
\[
[\ell_i^{(k)},u_i^{(k)}]
=
[\ell_i(5)-u_i(2),\;u_i(5)-\ell_i(2)].
\]
The load family has
\[
M_k=3\times4\times|\Bret|\times2=24|\Bret|.
\]

\textbf{Presentation order.} Define
\[
Q_i^{(o)}=D_i(o_Q)-D_i(o_A)\in[-2,2],
\]
with
\[
[\ell_i^{(o)},u_i^{(o)}]
=
[\ell_i(o_Q)-u_i(o_A),\;u_i(o_Q)-\ell_i(o_A)].
\]
The order family has
\[
M_o=3\times4\times|\Bret|\times2=24|\Bret|.
\]
Positive values mean DCI is larger when the added-slot occupant precedes the target. Negative values mean the reverse.

Every one of these families uses exact support $[-2,2]$ and width $w=4$. MP-EB therefore uses the $w=4$ range term and the family-specific $M$ in $\log(4M/0.05)$. WSR-PP uses normalized support $[-2,2]\mapsto[0,1]$ and per-tail allocation $0.05/(2M)$. The interval endpoints above are executable definitions and cannot be replaced by a narrower observed range.

\subsection{Direct paired order-effect family}\label{sec:orderfamily}
Order is a primary moderator in the primitive surface but the direct order difference is secondary-confirmatory. Its two-sided family is constructed from the paired $Q_i^{(o)}$ values above and is interpreted only after the primary headline gate $L_\Theta>0$. The gate cannot be bypassed by a large isolated order estimate. Order estimates and intervals are reported regardless of activation status, with their status labelled \texttt{DESCRIPTIVE\_UNTIL\_PRIMARY\_GATE} when the primary gate is not met.

\subsection{Non-claim-bearing bootstrap diagnostic}
A whole-superblock bootstrap may be run as an efficiency diagnostic to compare interval widths and empirical dependence structure. It resamples complete accepted superblocks within pair, preserving all references, budgets, loads, orders, and replicates. It is never substituted for the selected finite-sample confidence procedure, never rescues a failed claim, and never determines $H_p$ after confirmation.

\section{Prospective power and design simulation}\label{sec:power}

\subsection{Two simulation roles}
There are two distinct pre-confirmatory simulation exercises, both using 10,000 complete simulated experiments per scenario/design point and both calling the production analysis implementation.

\textbf{Inference-method comparison} occurs after the checkpoint/runtime and references are frozen but before construction of the full attainable superblock banks. It runs each statistically valid inference candidate through the same prespecified scenario set, family sizes, missingness logic, and planning $(\mathbf H,G)$ grid. Its purpose is to freeze one claim-bearing confidence method by the validity-first deterministic selection rule in Section~\ref{sec:analysis}.

\textbf{Authoritative design simulation} occurs after the selected method is frozen and the pair-specific attainable superblock banks establish maximum feasible counts. It searches only $(\mathbf H,G)$ designs that are actually attainable and selects the final pair-specific $H_p$ and common $G$. The inference method cannot be reopened at this stage because the attainable-bank outcome is observed later in the dependency graph.

\subsection{Frozen scenario dimensions}
Both exercises simulate the actual structured DCI experiment, rather than isolated normal means. The frozen scenario library includes at minimum: baseline target-failure probabilities $\{0.10,0.25,0.50,0.75,0.90\}$; headline effects $\Theta\in\{0,\epsilon,2\epsilon\}$; a single-reference bottleneck; heterogeneous primitive effects including a $1.5\times$ scale perturbation; shared focal/reference-arm dependence corresponding to correlation scenarios $\rho\in\{0,0.3,0.6\}$ where feasible; reference heterogeneity; the full retained budget/load/order grid; decoder replicate variation; 5\% boundary-contamination scenarios; and permanent technical-missingness rates $\{0,0.001,0.005,0.01,0.02\}$ under both outcome-independent and adverse outcome/length-linked mechanisms.

The simulator must generate focal-arm observations once and reuse them in the four paired focal-minus-reference primitives exactly as production analysis does. It must generate both loads and both orders within one superblock, propagate missing constituent generations into exact primitive intervals, preserve pair-specific $H_p$, and calculate every primary and secondary family with its real cardinality and support. The confirmatory background/task law is the independent with-replacement law in Section~\ref{sec:bgsampling}; no obsolete without-replacement cycle may survive in simulation code.

The central scenarios are defined prospectively in the simulation manifest. Stress scenarios are labelled before results are inspected. Simulation seeds, scenario IDs, code/environment hashes, and expected-compute model are frozen.

\subsection{Candidate $(\mathbf H,G)$ grids and compute accounting}\label{sec:hplanning}
The method-comparison stage uses a frozen finite planning grid wide enough to reveal whether either procedure requires tens, hundreds, or more superblocks, but it may not evaluate superblock counts already outside the plausible pre-Stage-10 construction domain. Before any Stage~9 method-comparison output is viewed, freeze for each pair a finite set $\mathcal C_p^{\rm pre}$ of hard construction ceilings already implied by pre-Stage-9 artifacts, including any finite content-inventory or duplicate-use ceiling, the predeclared Stage~10 accepted-bank construction cap, and any stricter proposal-space, runtime, storage, or cost ceiling established by the frozen construction/runtime plan. A category that supplies no finite hard ceiling is omitted. Define
\[
\widetilde H_{p,\max}^{\rm plan}
=
\min\mathcal C_p^{\rm pre}.
\]
The source artifact and integer value of every included ceiling, and the resulting $\widetilde H_{p,\max}^{\rm plan}$, are frozen before method-comparison outputs are viewed. This is an operational cap, not an estimated acceptance-rate model; no new statistical model is fit for attainability. Every Stage~9 planning design must satisfy
\[
H_p\le \widetilde H_{p,\max}^{\rm plan}
\qquad\text{for every pair }p.
\]
The exact finite $H$ and $G$ grid inside those caps is a \PFD{} artifact frozen from the same legitimate pre-Stage-9 information. It cannot be expanded, narrowed, or shifted after seeing which method wins, and Stage~10 realised attainable-bank outcomes are not used to define it.

The authoritative stage restricts candidate $H_p$ to values no larger than the frozen realised attainable count $H_{p,\max}$ for each pair; by construction $H_{p,\max}\le\widetilde H_{p,\max}^{\rm plan}$. Candidate $G$ values are a small frozen positive-integer grid justified by decoder Monte Carlo precision and compute. For every design, expected compute is calculated from measured valid logical generations/hour, empirical token counts from disjoint pre-confirmatory runtime material, retry/failure rate, and contemporaneous selected hardware cost. The scientific generation count remains
\[
N_{\rm conf}=20|\Bret|G\sum_{p=1}^{3}H_p.
\]
Cloud GPU count changes wall-clock parallelism, not the statistical $H_p$ or $G$.

\subsection{H-versus-G allocation is method-specific}
For MP-EB, the range penalty
\[
\frac{7w_r\log(2m/\alpha)}{3(H_r-1)}
\]
depends on independent superblocks and cannot be reduced by increasing $G$. Additional decoder replicates can reduce within-superblock Monte Carlo variation, but the design simulator must not spend compute on $G$ as though it were interchangeable with $H$ when the finite-sample range term dominates. The simulator therefore reports the variance-term and range-term contributions separately for representative primary primitives under every MP-EB candidate design.

For WSR-PP, there is no analytically assumed $H$-versus-$G$ exchange rate. The trade-off is determined by the same 10,000-experiment simulation because $G$ changes the distribution of the bounded superblock contrast while $H$ changes the number of capital-process factors. The method-comparison diagnostics also report, by scenario and design, the fraction of evaluated WSR-PP stake updates for which $\lambda_t^+(m)=c/m$ rather than $\widetilde\lambda_t$, including transition counts into and out of the truncation-bound regime. This diagnostic cannot be used to change $c$ or any other stake specification after comparison output is inspected. Planning figures such as $H\approx150$--$250$ have no preregistered status and cannot be cited as expected performance unless the frozen simulation actually produces them.

\subsection{Acceptance rule}
For a candidate method/design, the complete end-to-end simulation must satisfy all of the following using prespecified Monte Carlo uncertainty handling:
\begin{enumerate}
\item under $\Theta=0$, the simulated probability of the positive global claim $L_\Theta>0$ is no greater than 0.055;
\item under the central $\Theta=\epsilon$ scenarios, $P(L_\Theta>0)\ge0.80$;
\item under prespecified stress scenarios at $\Theta=\epsilon$, $P(L_\Theta>0)\ge0.70$;
\item under central $\Theta=2\epsilon$ scenarios, $P(L_\Theta>\epsilon)\ge0.80$;
\item all structural/inference invariants and separate-family multiplicity calculations remain valid under simulated missingness and dependence.
\end{enumerate}
The exact Monte Carlo acceptance convention at these thresholds, including whether a confidence interval around each simulated probability must lie on the passing side, is frozen before simulation. The same convention is used for both inference candidates and for authoritative design selection.

\subsection{Inference-method selection output}
For each statistically valid method, identify the lowest-expected-compute planning design within $H_p\le\widetilde H_{p,\max}^{\rm plan}$ that satisfies the acceptance rule. Compare those designs using the deterministic tie-breaking rule in Section~\ref{sec:analysis}. Archive, for every candidate method: validity status and reasons; power curves; median and upper-quantile interval widths; required planning $H_p$ and $G$; the frozen planning caps and their source artifacts; expected logical generations/tokens/GPU-hours/cost; WSR-PP truncation-bound frequencies where applicable; and sensitivity to missingness and dependence.

The chosen inference method is then frozen. No confirmatory outcome and no later attainable-bank shortfall can trigger method switching. If WSR-PP fails validity, MP-EB remains. If both are valid but only MP-EB passes, MP-EB remains. If both are valid and neither passes at feasible planning designs, MP-EB is retained as the conservative method and the downstream study can only become a feasibility-and-measurement pilot unless the later attainable bank unexpectedly admits a qualifying authoritative design within the already frozen planning domain.

\subsection{Pair-specific authoritative targets and floors}
After the attainable banks are built, the authoritative simulation evaluates only feasible $H_p\le H_{p,\max}$. The selected design is the lowest expected-compute valid design satisfying all acceptance criteria under the already frozen Stage~9 inference method. Pair-specific counts may differ. The complete selected target vector $\mathbf H=(H_1,H_2,H_3)$ and $G$ are frozen before the scientific-draw manifest is generated.

As a transparency diagnostic only, once Stage~10 has frozen the realised $H_{p,\max}$ values, reapply the already frozen deterministic Stage~9 method-selection criterion to the archived Stage~9 candidate-method simulation results after restricting the archive to designs satisfying $H_p\le H_{p,\max}$ for every pair. No new simulations, interpolation, extrapolation, tuning, candidate methods, or altered acceptance rules are permitted for this diagnostic. Record separately (i) the method actually frozen prospectively at Stage~9 and (ii) the valid method, if any, that would have won had the realised attainable region been known at Stage~9. If they differ, the final report must state the counterfactual method-selection discrepancy prominently. The discrepancy is disclosure only: it cannot reopen Stage~9, change claim-bearing inference, modify the final design to favour the counterfactual winner, or retune any candidate.

A pair-specific floor may be used only if it was defined as part of the frozen simulation candidate rule and the resulting unequal-$H_p$ family is simulated exactly. No pair is silently given fewer superblocks than the frozen final target because construction or runtime is inconvenient.

\subsection{If no design is adequately powered}
If no attainable design satisfies the acceptance rule under the selected valid method, the study proceeds only if frozen as a \textbf{feasibility-and-measurement pilot}. The SESOI is not reduced, the four-reference panels are not thinned, the search space is not narrowed based on expected results, multiplicity protection is not weakened, a less valid confidence method is not substituted, and additional $G$ is not used rhetorically to compensate for inadequate $H$. The report must state prospectively that the design did not meet the positive-headline power criteria.

\subsection{Required simulation outputs}
The freeze bundle includes machine-readable scenario definitions; all 10,000-replicate summary tables; null false-positive estimates; power at $\epsilon$ and $2\epsilon$; interval-width distributions; method-specific $H$-versus-$G$ curves; the frozen $\widetilde H_{p,\max}^{\rm plan}$ values and source ceilings; MP-EB range/variance decomposition where applicable; WSR-PP truncation-bound frequencies and transition diagnostics where applicable; missingness sensitivity; pair-specific realised $H_{p,\max}$ overlays on the already archived Stage~9 candidate curves; selected method/design IDs; expected compute; the prospective Stage~9 method; the counterfactual winner obtained by reapplying the frozen method-selection criterion within $H_p\le H_{p,\max}$; an explicit same/different indicator; and the full code/environment hash chain. Any counterfactual discrepancy is reported prominently but never changes the frozen method or design-selection rules. These are pre-confirmatory design artifacts, not evidence about realised DCI.

\section{Ordered construction, freeze, confirmation, and falsification pipeline}\label{sec:stages}
The ordering below is mandatory. A later stage may consume frozen outputs from earlier stages; it may not retrospectively change them. Any change to checkpoint, tokenizer, chat template, decoder, controller semantics, $V$, retained budgets, accelerator backend, GPU class, precision, attention backend, or other behaviourally relevant runtime field invalidates every downstream model-dependent artifact identified in the dependency manifest.

\subsection{Stage 1: scientific definitions, registries, splits, SESOI floor, and selection criteria}
Freeze the three target--focal rule definitions, exactly-four-reference requirement, family registries/schema mappings where already known, candidate/reserve authoring rubrics, background strata, material splits, field-role registry, proposal distributions, decoder specification, SESOI floor $\epsilon\ge0.06$, model/checkpoint eligibility criteria, Lambda hardware benchmark rule, inference-candidate specifications, and terminal-falsification design template. Establish content sufficiency for every scientific split. No model-dependent DCI outcome is generated for design selection.

\subsection{Stage 2: satisfiability and verifier infrastructure}
Implement/reconcile all deterministic verifiers, structured feature encodings, Z3 conjunctions where applicable, automatic witness generation, verifier conjunction checks, semantic-review triggers, boundary/property tests, causal field roles, completeness ledger, and provenance. Pair~1 bullet-count/transitive-order support must be built; Pair~2 and Pair~3 existing verifier code may be reused only after hash/semantic audit. Freeze the exact inherited binomial reliability implementation or resolve it as a pre-freeze output before dependent reliability audits.

\subsection{Stage 2A: early scientific-runtime and inference qualification}
Before expensive model-dependent construction, qualify candidate checkpoints, tokenizer/chat-template fixtures, decoder, reasoning/final-answer controller, candidate budget envelopes, $V$ procedure, Transformers reference engine, candidate accelerator backends, and both inference implementations. Run, wherever technically possible, all tests that could invalidate later calibration: tokenizer/template fixtures; reasoning-cap $b-1,b,b+1$ boundaries; final-answer $V-1,V,V+1$ boundaries; parser/controller property tests; same-device replay; cross-device replay/interchangeability on candidate hardware; maximum-context/OOM tests; dtype and attention-backend tests; deterministic-runtime flags; manifest fail-closed mutations; worker-death and identical-draw recovery; durable output persistence; verifier replay from stored bytes; seed collision tests; interval-missingness brute force; zero-variance inference fixtures; WSR-PP truncation-transition fixtures; and max--min small-family fixtures.

At this stage a betting method is only a \PROV{} candidate. Passing unit fixtures does not yet select it; the full validity mapping and prospective comparison occur later. Any candidate inference implementation that fails basic mathematical/software invariants is removed from eligibility before downstream work.

\subsection{Stage 3: outcome-independent checkpoint and accelerator/runtime freeze}
Select the confirmatory checkpoint using the frozen non-DCI eligibility criteria. Execute the Lambda hardware benchmark in Section~\ref{sec:hardwareselection} on the exact maximum scientific envelope. Select and freeze the lowest-cost eligible configuration. Certify the intended GPU pool as one interchangeability class under Section~\ref{sec:interchange}. Freeze the exact accelerator backend, GPU SKU, driver/CUDA or ROCm stack, PyTorch, Transformers, attention backend, dtype, container/dependency hashes, model/tokenizer/template, decoder, controller, verifier stack, seed mapper, and persistent-storage layout.

If no candidate hardware/runtime passes correctness, maximum-context, and interchangeability gates, stop before model-dependent scientific calibration. Do not fall back to a non-interchangeable large confirmatory pool.

\subsection{Stage 4: budget manipulation and final-answer allowance}
On disjoint pre-confirmatory material, freeze $V$ and run the manipulation procedure over the authored candidate hard-budget grid. Retain the largest ordered $\Bret$ satisfying the frozen adjacent-separation criterion with at least three budgets. Failure to retain at least three budgets makes the checkpoint/runtime ineligible and requires returning to Stage~3 under the already frozen checkpoint-selection rule, before reference construction begins.

\subsection{Stage 5: exact SESOI if required, then boundary parameter calibration}
If exact $\epsilon$ has not yet been scientifically frozen, it may remain unresolved only while no numerical procedure consumes it. Before the first exact dependency, execute and sign the prospective scientific-decision procedure in Section~\ref{sec:sesoi}, enter the exact value and rationale in the decision record, and then freeze exact $\epsilon$. Run the fixed $n_{\rm cal}=100$ Clopper--Pearson boundary procedure for target, focal, background, and candidate-reference parameter grids across relevant budgets/orders. Freeze nonempty family-specific $P_Q$ and $\nu_Q$ outputs.

\subsection{Stage 6: candidate-reference hard eligibility and coupling}
Start with 24 authored candidates per pair. Apply logical, schema, scope, satisfiability, boundary, and verifier gates. Execute the fixed coupling probe with 2,000 base outputs and eight nuisance classes. Because $\epsilon\ge0.06$, the coupling acceptance ceiling is already $0.015$. If the fixed probe-expansion branch requires the $\epsilon/8$ precision target, exact $\epsilon$ must have been frozen before any expansion data are generated. Expansion uses the frozen smallest-multiple-of-500 rule. No threshold changes to force eligibility.

\subsection{Stage 7: semantic paraphrases, wording instability, and candidate-pool count}
Certify one canonical wording plus 11 semantic paraphrases for every surviving relevant constraint. Estimate wording instability with $N_W=100$, $G_W=2$, whole-task cluster resampling, and 10,000 bootstrap replicates. Since $\epsilon\ge0.06$, the fixed wording-difference ceiling is $0.02$. Require at least 12 hard-eligible candidate references per pair; if fewer survive from the first 24, add exactly one blind tranche of 12 and repeat the same hard gates. Fewer than 12 of 36 emits \REFPOOL{}.

\subsection{Stage 8: reference-matching calibration, simultaneous admission, and independent validation}
Freeze exact $\epsilon$ before any matching caliper, sensitivity coefficient, or joint reference-error allowance whose numerical value depends on it. Estimate the four matching coordinates using complete construction clusters $U_i$ and 10,000 cluster bootstraps. Apply simultaneous uncertainty-aware calipers, joint reference-error certificate, deterministic ranking among already eligible candidates, and exact tie rule. Admit exactly four references per pair. Run the complete disjoint independent-validation procedure. Failure returns the affected pair to the already frozen pre-confirmatory construction process; confirmatory outcomes still do not exist.

\subsection{Stage 9: inference-method qualification and selection}
Complete the theorem/assumption mapping for MP-EB and WSR-PP under the now-frozen superblock law, support definitions, decoder, reference family, and runtime. Exact $\epsilon$ must already have been fixed by the procedure in Section~\ref{sec:sesoi}. Before any comparison output is viewed, freeze the pair-specific $\widetilde H_{p,\max}^{\rm plan}$ caps and their source artifacts under Section~\ref{sec:hplanning}. Execute all validity fixtures, including the validity-critical WSR-PP truncation-transition fixtures, and the 10,000-experiment candidate-method comparison using the same frozen scenarios and planning $(\mathbf H,G)$ grid restricted to $H_p\le\widetilde H_{p,\max}^{\rm plan}$. Select the method by the deterministic validity-first rule. Freeze formulas, betting stake/truncation constants if WSR-PP wins, MP constants if MP-EB wins, numerical inversion, multiplicity allocation, interval-completion implementation, primitive ordering, production/reference code hashes, and method-selection report.

No later attainable-bank result may cause method re-selection. The post-Stage~10 restricted-region method comparison is a mandatory transparency diagnostic only. If no betting candidate validates, retain MP-EB. If no candidate validates, stop.

\subsection{Stage 10: attainable superblock-bank construction}
Under the independent with-replacement proposal law, construct a reserve bank for each pair using domain-separated proposal streams and the full whole-composition acceptance gate. Generate and verify concrete satisfying witnesses automatically where deterministic coverage permits. Archive every proposal and rejection. Continue only to the predeclared construction cap already incorporated into $\widetilde H_{p,\max}^{\rm plan}$ and determine each pair's realised attainable count $H_{p,\max}\le\widetilde H_{p,\max}^{\rm plan}$. Do not change proposal probabilities, pool composition, acceptance rules, construction cap, or Stage~9 method in response to the number or identity of successful superblocks.

\subsection{Stage 11: authoritative end-to-end $(\mathbf H,G)$ design simulation}
Run 10,000 simulated complete experiments per attainable candidate design through the single frozen inference method and exact production analysis. Use actual family cardinalities, shared focal dependence, background/task law, interval missingness, max--min headline, and measured runtime-cost model. Explicitly evaluate compute allocation between independent superblocks and decoder replicates. Before selecting the authoritative design, compute the mandatory counterfactual method-selection diagnostic by restricting the archived Stage~9 candidate curves to $H_p\le H_{p,\max}$ and reapplying the frozen Stage~9 criterion without rerunning or retuning either method. Archive and report any discrepancy, but do not reopen method selection. Select the lowest expected-compute valid design satisfying every prospective acceptance criterion under the method actually frozen at Stage~9. If none exists, freeze the feasibility-and-measurement-pilot interpretation.

\subsection{Stage 12: freeze numerical pre-confirmatory outputs and scientific-draw manifest}
Freeze final $H_p$, $G$, exact $\epsilon$, $P_Q$, $\nu_Q$, four reference identities per pair, matching constants, retained budgets, $V$, selected inference method, runtime/hardware identities, terminal-falsification fixtures/rules, and every other pre-confirmatory numerical artifact. Select exactly the first $H_p$ eligible accepted superblocks from each pair's bank under the frozen bank-order rule. Expand them into the complete generation manifest with immutable IDs and seeds. Compute the root manifest hash.

\subsection{Stage 13: final runtime re-certification}
Re-run the complete mandatory runtime and inference-validation suite on the already qualified frozen scientific configuration and intended confirmatory GPU pool. Stage~13 is re-certification, not the first correctness test. Compare all hashes to Stage~2A/3. Re-run maximum-context, same-device/cross-device replay, interchangeability, worker-death recovery, persistence, seed, controller, verifier replay, inference fixtures, and manifest fail-closed tests. Any failure blocks confirmation and invalidates the dependency-scoped downstream artifacts; it does not license silent repair after generation begins.

\subsection{Stage 14: immutable protocol and analysis freeze}
Sign the preregistration source/PDF, complete decision register, remaining-pre-freeze closure report, construction artifacts, runtime manifest, inference code, simulation outputs, scientific-draw manifest, terminal-falsification plan, and root SHA-256 manifest. Record that no confirmatory output exists before the signed freeze timestamp.

\subsection{Stage 15: confirmatory generation}
Execute only the frozen generation manifest. The scheduler may distribute independent generation IDs across certified GPUs but cannot modify scientific IDs, seeds, prompts, runtime, model, batching policy, or stopping. Worker/device failures use identical-draw replay within the certified interchangeability class. Technical missingness is assigned only after the frozen retry policy is exhausted.

\subsection{Stage 16: locked confirmatory inference}
Verify completeness and integrity, freeze the terminal observed/missing ledger, and run the single preselected analysis code once as the authoritative confirmatory analysis. Produce all primary and secondary family bounds, the full reference-specific surface, least-favourable bounds, $L_\Theta$, selected descriptive maximum, result/status codes, and prespecified sensitivities. The confirmatory analysis cannot trigger new generation, method selection, or sample-size extension.

\subsection{Stage 17: sequestered terminal falsification}
Only after the confirmatory outputs and analysis hash are locked may the terminal-falsification material be opened. Execute the frozen battery and decision rules. Terminal falsification can qualify, challenge, or prevent scientific interpretation under its prespecified labels, but it cannot trigger selective replacement of confirmatory data or retroactive redesign. Any concern is reported alongside the locked confirmatory result under \TERMFALSE{} or its frozen subcodes.

\section{Stopping, fail-closed behaviour, and non-identification}\label{sec:failclosed}

\subsection{Fixed-size stopping}
For each pair, confirmation ends when every generation in the frozen manifest corresponding to exactly $H_p$ selected superblocks and $G$ scheduled replicates has either one durably persisted scientific output or a terminal permanent-technical-missingness record after the frozen retry policy. There is no optional stopping on effect size, significance, confidence width, favourable reference, runtime convenience, or apparent stability.

\subsection{Run-level integrity failure}
Emit \RUNFAIL{} when the authoritative scientific manifest cannot be reconstructed or trusted, including an unapproved change in checkpoint/runtime/decoder/controller, corruption or unexplained loss of scientific identities, unresolvable arm mismatch, generation under an uncertified device/runtime, or evidence that outputs were selectively regenerated. A run-level integrity failure blocks confirmatory scientific claims until a new pre-confirmatory protocol and run are established.

\subsection{Primary inferential failure}
\INFERNE{} is reserved for structural conditions that make the frozen claim-bearing procedure undefined or its required theorem assumptions false. Examples include discovering after freeze that the implemented superblock sampler is not the frozen independent with-replacement law; using an inference implementation whose frozen validity artifact is corrupted or not the selected hash; $H_p<2$ where the selected method requires at least two observations; or an unresolvable family-cardinality/support mismatch.

Zero empirical variance is not an inferential failure. Permanent technical missingness is not an inferential failure merely because intervals widen. A valid but uninformative lower bound remains an estimable negative or wide result. No post-confirmatory method switch is permitted as a remedy for weak precision.

\subsection{Technical-integrity warning}
Emit \TECHWARN{} when technical missingness or retry behaviour exceeds the frozen warning threshold while the run remains analysable through interval completion. Report warning rates by pair, budget, load, order, arm, device, and failure category. The warning does not itself delete data or change the confidence method.

\subsection{Region-level non-identification}
A cell or pair receives \NOTID{} only for a pre-confirmatory construction failure such as inadequate reference pool, no valid four-reference panel, insufficient satisfiable material, failed verifier/reliability gate, or no valid retained budget treatment. Such regions never enter $\Sc$. They cannot be dropped after confirmation merely because their effects are inconvenient.

\subsection{Arm mismatch}
Any divergence in a pretreatment field required to be common between a focal and its paired reference arm emits \ARMMISMATCH{}. The mismatch is an execution defect, not a covariate imbalance. The affected scope follows the fail-closed integrity policy and is never statistically adjusted away.

\section{Result taxonomy and reporting logic}

\subsection{Primary headline states}
After locked inference, label the global result as exactly one of:
\begin{itemize}
\item \texttt{ROBUST\_POSITIVE\_SESOI}: $L_\Theta>\epsilon$;
\item \texttt{ROBUST\_POSITIVE}: $0<L_\Theta\le\epsilon$;
\item \texttt{NOT\_ESTABLISHED}: $L_\Theta\le0$ with an otherwise valid confirmatory run;
\item \texttt{FEASIBILITY\_MEASUREMENT\_PILOT}: the pre-confirmatory authoritative design failed the prospective power rule, regardless of descriptive confirmatory magnitude;
\item \INFERNE{} or \RUNFAIL{} when the corresponding structural failure applies.
\end{itemize}
The labels are mutually exclusive at the headline level. A technical-integrity warning may coexist with a scientifically estimable state.

\subsection{Cell-level states}
For every retained cell, report all four reference-specific estimates and simultaneous lower bounds, the least-favourable point summary, and the least-favourable lower bound. A cell is reference-robust positive only when all four lower bounds exceed zero; it exceeds the SESOI only when all four exceed $\epsilon$. Mixed signs are explicitly labelled reference-dependent and are never averaged into a primary control.

\subsection{Order, compute, and load states}
Report the separately protected two-sided primitive, compute, load, and order intervals under their own familywise guarantees. State prominently that these families each control their own FWER at 0.05 and that the collection of all primary and secondary families does not have a single joint 5\% FWER. Direct order claims are labelled confirmatory only when the primary $L_\Theta>0$ gate is active; otherwise their intervals are reported as prespecified but hierarchically inactive.

\subsection{Apparatus and identification states}
Report reference-pool insufficiency, verifier/reliability failure, satisfiability failure, controller/runtime failure, interchangeability failure, vLLM equivalence failure, excessive technical missingness, and terminal-falsification concern as apparatus/identification states distinct from behavioural DCI outcomes. Optional-vLLM failure simply leaves Transformers as the scientific engine and is not a scientific failure.

\subsection{Required selection-bias adjacency rule}
Whenever $\widehat\Theta_{\rm crit}$ or its selected cell $s^\ast$ appears, the same paragraph/table/figure panel must display $L_\Theta$ and the statement that the maximised plug-in estimate is selected from a noisy surface and can overstate the underlying maximised effect. The claim-bearing statement is based on the simultaneous lower-bound transformation, not on the magnitude of $\widehat\Theta_{\rm crit}$.

\section{Pre-freeze decision register}\label{sec:register}
\small
\begin{longtable}{p{0.04\textwidth}p{0.27\textwidth}p{0.19\textwidth}p{0.12\textwidth}p{0.285\textwidth}}
\toprule
\textbf{ID} & \textbf{Decision/output} & \textbf{Status} & \textbf{Freeze blocker?} & \textbf{Rule} \\
\midrule
\endhead
D1 & DCI replacement estimand and failure scale & \FIXED{} scientific commitment & No & $A+B+S$ versus $A+R_j+S$; positive means more target failure. \\
D2 & Three target--focal pairs & \FIXED{} scientific commitment & No & Exact bullet count/transitive order; mandatory words/all-words-unique; sentence count/increasing lengths. \\
D3 & Reference architecture & \FIXED{} scientific commitment & No & Exactly four independently admitted references; primitive effects remain reference-specific; no primary averaging/weights. \\
D4 & Robust/headline operators & \FIXED{} scientific commitment & No & Least-favourable reference within cell; global maximum over frozen cells. \\
D5 & Loads and orders & \FIXED{} scientific commitment & No & Paired $k\in\{2,5\}$ and both target/added-slot orders within superblock. \\
D6 & Constraint parameters & \FIXED{} method; numerical outputs pending & Yes & Family-specific $P_Q,\nu_Q$ from boundary calibration; paired realisations. \\
D7 & Universal cross-family severity coordinate & \FIXED{} excluded & No & Do not match heterogeneous parameter severity as a scalar coordinate. \\
D8 & Candidate-reference pool & \FIXED{} method & No & 24 plus at most one blind tranche of 12; require at least 12 hard-eligible before matching. \\
D9 & Construction bootstrap & \FIXED{} method & No & Complete $U_i$ clusters; exactly 10,000 replicates. \\
D10 & Verifier coupling & \FIXED{} method; realised $n$ pending & Yes if expansion occurs & Eight nuisance classes; start 2,000; ceiling $0.015$ because $\epsilon\ge0.06$; exact $\epsilon$ before any $\epsilon/8$ expansion. \\
D11 & Wording instability & \FIXED{} method & No & $M=12$, $N_W=100$, $G_W=2$, covariance-corrected estimate, 10,000 task-cluster bootstraps; ceiling $0.02$. \\
D12 & Reference matching/admission/validation & \FIXED{} method; constants/identities pending & Yes & Frozen coordinates, calipers, certificate, deterministic ranking, exactly four, disjoint validation. \\
D13 & Whole-composition satisfiability & \FIXED{} method; artifacts pending & Yes & Abstract solver evidence where applicable + concrete all-verifier-pass witness + semantic review only where deterministic coverage is incomplete. \\
D14 & Confirmatory superblock probability law & \FIXED{} method & No & Independent with-replacement task/background proposal streams; whole-proposal rejection sampling; no round-robin WoR cycles. \\
D15 & Scientific unit identity & \FIXED{} method & No & Accepted draw index makes repeated contexts distinct independent superblocks. \\
D16 & Confirmatory checkpoint/revision & \PFD{} & Yes & Select only by frozen non-DCI capability, licensing, fit, reproducibility, and resource criteria. \\
D17 & Scientific software architecture & \FIXED{} method & No & PyTorch + Hugging Face Transformers + frozen accelerator backend. \\
D18 & Planned scalable substrate & \FIXED{} plan & No & Lambda On-Demand Cloud; exact availability/prices reverified immediately before freeze. \\
D19 & Accelerator/GPU/runtime manifest & \PFD{} / \GATE{} & Yes & Outcome-independent benchmark; freeze backend, GPU SKU, driver, CUDA/ROCm, PyTorch, Transformers, attention, dtype, container, model, tokenizer, template, controller, decoder, and verifiers. \\
D20 & GPU interchangeability class & \GATE{} & Yes & Exact cross-device replay on intended pool; no non-interchangeable physical-GPU pinning fallback. \\
D21 & Tensor parallelism & \FIXED{} default excluded & Yes if invoked & Prefer single-GPU checkpoint; TP requires separate pre-freeze necessity justification and complete validation. \\
D22 & Optional vLLM & \PFD{} / \GATE{} & No & Only if full runtime suite + 10,000 paired prompt-seed exact equivalence passes; otherwise Transformers remains. \\
D23 & Stochastic decoder & \FIXED{} method & No & $T=.70$, top-$p=.95$, top-$k$ disabled, one beam, no extra truncation/selection processors. \\
D24 & Hard reasoning controller and $V$ & \FIXED{} controller method; $V$ pending & Yes & Forced transition at $b$, fixed final-answer allowance, no cross-generation state leakage. \\
D25 & Retained model-specific budget grid & \PFD{} / \GATE{} & Yes & At least three conditions pass the frozen manipulation gate. \\
D26 & SESOI floor & \FIXED{} scientific commitment & No & $\epsilon\ge0.06$ from Stage~1. \\
D27 & Exact SESOI & \PFD{} until first exact dependency & Yes & Select only by the signed prospective scientific-decision procedure in Section~\ref{sec:sesoi}; freeze rationale and value before $\epsilon/8$ expansion, any exact-$\epsilon$ construction dependency, inference comparison, or authoritative design selection. \\
D28 & Confirmatory inference candidates & \FIXED{} candidate specifications & No & MP-EB and WSR-PP only. \\
D29 & Selected claim-bearing inference method & \PFD{} / \GATE{} & Yes & Validity first; then 10,000-experiment common comparison within frozen pre-Stage-10 $H$ caps; freeze method/formulas/code hash before full bank. Later realised attainability cannot reselect the method; restricted-region counterfactual discrepancy is disclosure only. \\
D30 & Zero-variance handling & \FIXED{} method & No & Ordinary case for both candidates; no studentised denominator. \\
D31 & Permanent technical missingness & \FIXED{} method & No & Identical-draw retry then support-valid interval completion; no deletion or replacement seed; no parametric MCAR/MAR or contemporaneous outcome-independence assumption, conditional on the frozen latent complete-data filtration/conditional-mean property. \\
D32 & Exact primitive support & \FIXED{} method & No & $D\in[-1,1]$, width 2. \\
D33 & Exact compute/load/order supports & \FIXED{} method & No & Each $Q\in[-2,2]$, width 4; explicit interval endpoints in Section~\ref{sec:secondarysupports}. \\
D34 & Multiplicity policy & \FIXED{} method & No & Primary family own 5\% FWER; each secondary-confirmatory family separately 5\%; no joint all-family 5\% claim; order hierarchically activated. \\
D35 & Inference-method comparison & \FIXED{} method; output pending & Yes & 10,000 experiments/scenario/design, same scenarios/family sizes, validity before efficiency; planning grid obeys frozen $H_p\le\widetilde H_{p,\max}^{\rm plan}$ and archives later restricted-region counterfactual winner. \\
D36 & Pair-specific $H_p$ and common $G$ & \PFD{} & Yes & Select only after attainable banks with 10,000-experiment authoritative simulation. \\
D37 & H-versus-G allocation & \FIXED{} method & No & MP range term cannot be reduced by $G$; WSR trade-off determined empirically by exact simulation. \\
D38 & Feasibility-pilot fallback & \FIXED{} method & No & If no attainable valid design passes, do not weaken SESOI, references, multiplicity, search space, or validity. \\
D39 & Selected maximum reporting & \FIXED{} method & No & $\widehat\Theta_{\rm crit}$ descriptive only; $L_\Theta$ claim-bearing and adjacent. \\
D40 & Runtime validation timing & \FIXED{} method & No & Invalidating tests at Stage~2A/3; Stage~13 final re-certification. \\
D41 & Failure recovery/storage & \FIXED{} method & No & Replay identical draw on certified interchangeable device; persistent scientific ledger/raw outputs/hashes; ephemeral cache only. \\
D42 & Final scientific-draw manifest/root hash & \PFD{} & Yes & Generated only after final $H_p,G$ and all numerical outputs freeze. \\
D43 & Terminal-falsification battery & \PFD{} artifact; \FIXED{} sequencing & Yes & Freeze and sequester before confirmation; open only after confirmatory inference is locked. \\
D44 & Implementation readiness & engineering status & Yes through gates & Reusable provenance/construction core; Pair~2/3 verifiers present; Pair~1, inference layer, controller, and superblock spine require completion/validation. \\
\bottomrule
\end{longtable}
\normalsize

The register distinguishes scientific commitments from fixed methods, deterministic selection procedures, numerical/empirical outputs, validation gates, secondary analyses, and future work. A row marked \PFD{} is not permission for implementer discretion: its value is produced only by the rule cited in that row.

\section{Fixed, provisional, secondary, unresolved, and future components}

\subsection{Fixed scientific commitments}
Fixed commitments include: causal replacement rather than deletion; the three target--focal pairs; exactly four separate references; no primary reference weighting; least-favourable-reference robustness; global max--min headline; paired loads and both orders; variable family-specific parameters; exclusion of universal severity matching; one pinned open-weight checkpoint for confirmation; model-specific hard reasoning budgets; fixed $V$ across budgets; stochastic decoder; no best-of-$N$/reranking/revision; behavioural/technical failure separation; fixed-size stopping; immutable scientific identities; whole-composition satisfiability; pair identity descriptive; SESOI floor $\epsilon\ge0.06$; and feasibility-pilot fallback if prospective power fails.

\subsection{Fixed construction and analysis methods}
Fixed methods include the reference-construction firewall, eligibility/candidate rules, construction-cluster bootstrap, coupling audit, wording-instability estimator, independent validation, independent with-replacement confirmatory superblock law, interval completion, exact supports, separate-family multiplicity policy, max--min bound transformation, selected-maximum warning, hardware benchmark rule, interchangeability/retry design, and runtime-validation sequencing.

\subsection{Deterministic pre-freeze selection procedures}
The prospective SESOI documentation rule and the procedures for checkpoint choice, $V$ and budget retention, hardware/runtime choice, pre-Stage-10 $H$ planning caps, inference-method qualification/selection, reference ranking, final $H_p,G$ choice, and optional vLLM substitution are fixed before their dependent outputs are viewed. Their realised outputs remain unknown until the required pre-confirmatory evidence is generated.

\subsection{Genuinely remaining numerical, artifact, and gate outputs}
The unresolved outputs are listed in the accompanying \nolinkurl{DCI_Pilot_Preregistration_v19_REMAINING_PREFREEZE.md}. They include exact registry mappings where unresolved; calibrated $P_Q,\nu_Q$; numerical matching constants; four actual references and validation outcomes; checkpoint/revision; exact $V$ and $\Bret$; exact $\epsilon$ and signed SESOI rationale if not yet frozen; hardware/backend/runtime manifest and interchangeability outcome; pre-Stage-10 $\widetilde H_{p,\max}^{\rm plan}$ ceilings; inference-method validity results and selected method/hash; any coupling-bank expansion; attainable banks and $H_{p,\max}$; the realised-attainability counterfactual method-selection diagnostic; final $H_p,G$; optional-vLLM outcome; final scientific-draw manifest; and sequestered terminal-falsification artifacts.

\subsection{Secondary analyses}
Prespecified secondary work includes the separately protected two-sided primitive, compute, load, and direct order families; adjacent-budget descriptive contrasts unless a separate family is frozen before confirmation; focal/reference compliance; behavioural state decompositions; realised reasoning, cap-hit, answer-length, latency, and technical-missingness diagnostics; whole-superblock bootstrap efficiency diagnostics; descriptive cross-pair comparisons; and clearly labelled exploratory mediation or hierarchical models. None can alter the primary search set, selected confidence method, or headline.

\subsection{Future work}
Future versions may add more loads, independently sample target--focal pairs from constraint-family populations, test reverse-direction DCI, transport across model families under separate budget calibration, investigate mechanisms, evaluate safety-monitor prompts, or develop stronger finite-sample procedures. Such work is outside v19 and cannot be smuggled into the current confirmation after outcomes are visible.

\section{Claim boundaries and escalation criteria}

\subsection{Strongest supported interpretation}
A successful positive result supports a narrow causal statement: for the frozen checkpoint, runtime/backend, stochastic decoder, accepted-superblock population, hard reasoning cap, load, and presentation order, replacing each of the four independently admitted references with the focal constraint increases target-failure probability under the prespecified reference-specific causal contrast. The simultaneous confidence guarantee is exactly the guarantee of the selected primary family, followed by the deterministic least-favourable and max--min transformations.

If $L_\Theta>0$, at least one preregistered cell supports reference-robust positive DCI at the primary familywise level. If $L_\Theta>\epsilon$, the corresponding lower bound exceeds the preregistered SESOI. The selected descriptive maximum is not substituted for either claim.

\subsection{What the pilot cannot establish}
The pilot cannot establish: a universal law of constraint interaction; population-level effects of target or focal families; transport to untested pairs/tasks/loads/wordings/background populations; transport of raw reasoning-token counts across tokenizers or model families; monotone compute effects; an internal mechanism; mediation by realised reasoning or answer length from the primary estimand; reverse-direction interference unless separately preregistered; behaviour under unsatisfiable compositions; evidence of absence from a bound that overlaps zero; a single 5\% FWER jointly across all primary and secondary families; or that a reference-dependent profile has one scientifically privileged average.

A betting method being selected cannot establish that betting intervals are generally superior. It establishes only that the prespecified candidate passed the v19 validity gates and won the frozen prospective comparison for this design. Conversely, retaining MP-EB establishes no defect in betting methods generally.

\subsection{Escalation from local baseline to confirmatory checkpoint and Lambda execution}
Local open-weight baselines are engineering qualification tools. A checkpoint/runtime may proceed to confirmatory construction only after: reasoning/final-answer separation is deterministic; hard-cap accounting passes; at least three budget conditions can be retained; verifier and answer extraction are reliable; the maximum prompt/budget/$V$ envelope fits candidate hardware; the selected accelerator/runtime passes Stage~2A tests; and the outcome-independent Lambda benchmark identifies an eligible confirmatory configuration. Local DCI magnitude is sequestered from checkpoint, reference, inference-method, SESOI, $H_p$, $G$, and hardware decisions.

The scalable execution step changes capacity, not the scientific protocol. Horizontal workers process immutable independent generation IDs under one frozen interchangeability class. GPU count may be adjusted before the first confirmatory generation only within the certified class and without changing model process semantics, batching, dtype, or other behaviourally relevant settings. After confirmation starts, adding a device requires that it already belongs to the frozen certified class or pass the prespecified replacement-device certification without changing the class manifest.

\section{Compact pre-freeze checklist}

\subsection{Scientific design and dependency order}
\begin{itemize}
\item[$\square$] Three fixed target--focal pairs, exactly four references per pair, both orders, and paired loads $k=2,5$ are frozen.
\item[$\square$] All material splits and content-sufficiency counts are frozen and disjoint.
\item[$\square$] SESOI floor $\epsilon\ge0.06$ is signed at Stage~1; exact $\epsilon$ and its signed prospective scientific rationale were frozen under Section~\ref{sec:sesoi} before the first exact numerical dependency.
\item[$\square$] Universal cross-family parameter-severity matching is absent.
\item[$\square$] Confirmatory superblock proposals use the frozen independent with-replacement law and whole-proposal rejection sampling.
\item[$\square$] The signed dependency manifest shows that no downstream artifact predates a behaviourally relevant upstream runtime/model change.
\end{itemize}

\subsection{Reference construction, verifiers, and satisfiability}
\begin{itemize}
\item[$\square$] Pair~1 bullet-count/transitive-order verifiers and all reused Pair~2/3 verifiers pass semantic, boundary, differential, and property tests.
\item[$\square$] Candidate pool uses 24 plus at most one blind tranche of 12 and has at least 12 hard-eligible survivors before matching.
\item[$\square$] Construction uncertainty resamples complete $U_i$ clusters with exactly 10,000 replicates.
\item[$\square$] Coupling audit passes all error and $\kappa$ ceilings; any probe expansion occurred only after exact $\epsilon$ freeze.
\item[$\square$] Twelve semantic wordings per relevant constraint and $100\times2$ wording measurements pass the frozen $0.02$ ceiling.
\item[$\square$] Calipers, sensitivity coefficients, joint certificate, deterministic ranking, exactly-four selection, and independent validation are frozen.
\item[$\square$] Every confirmatory composition has required solver evidence, a concrete all-verifier-pass witness, and human semantic review only where deterministic coverage is incomplete.
\end{itemize}

\subsection{Decoder, runtime, Lambda, and recovery}
\begin{itemize}
\item[$\square$] Exact checkpoint/revision, tokenizer, chat template, PyTorch, Transformers, accelerator backend, GPU SKU, driver, CUDA/ROCm version, attention backend, dtype, controller, parser, verifiers, container, and environment flags are frozen.
\item[$\square$] Lambda hardware candidates passed maximum-context and complete correctness/reproducibility tests; lowest-cost eligible configuration selected by the frozen rule using archived contemporaneous prices.
\item[$\square$] Intended confirmatory GPUs form one certified interchangeability class with exact cross-device replay; no device-pinning fallback remains.
\item[$\square$] Persistent storage contains manifests, raw outputs, attempt lineage, completed ledger, verifier outputs, and hashes; ephemeral storage contains only reproducible scratch/cache.
\item[$\square$] Decoder is exactly $T=.70$, top-$p=.95$, top-$k$ disabled, one beam, with no extra truncation/selection processors.
\item[$\square$] SHA-256 seed map and canonical serialization are frozen and collision-tested.
\item[$\square$] Reasoning and final-answer boundary tests, same/cross-device replay, OOM, dtype/attention, worker-death recovery, manifest fail-closed, output persistence, and verifier replay all pass at Stage~2A/3 and again at final re-certification.
\item[$\square$] vLLM is absent or has passed the full suite plus 10,000 exact prompt-seed equivalence cases before model-dependent construction.
\end{itemize}

\subsection{Statistical integrity}
\begin{itemize}
\item[$\square$] MP-EB and WSR-PP have complete theorem/assumption mappings to the actual superblock law.
\item[$\square$] Selected inference method passed exact-support, zero-variance, interval-missingness, familywise multiplicity, max--min, null-coverage, and independent-reference-implementation gates; WSR-PP, if eligible, also passed all validity-critical truncation-transition fixtures.
\item[$\square$] Method-selection comparison used exactly 10,000 simulated experiments per frozen scenario/design point, obeyed the prospectively frozen $H_p\le\widetilde H_{p,\max}^{\rm plan}$ caps, and was completed before attainable-bank authoritative design selection.
\item[$\square$] Primary contributions use support $[-1,1]$; compute/load/order contributions use $[-2,2]$ and the exact interval-completion formulas in Section~\ref{sec:secondarysupports}.
\item[$\square$] Primary family controls its own 5\% FWER; each secondary-confirmatory family separately controls its own 5\% FWER; no joint-all-families 5\% statement appears.
\item[$\square$] Permanent technical missingness uses interval completion; no superblock deletion, replacement seed, parametric MCAR/MAR model, contemporaneous outcome-independence assumption, or complete-case headline analysis; the frozen latent complete-data conditional-mean property with respect to the full past filtration is verified.
\item[$\square$] Authoritative 10,000-experiment simulation used the Stage~9 selected method, actual family sizes, shared focal dependence, missingness, max--min, and actual attainable $H_{p,\max}$; the archived Stage~9 curves were also restricted to the realised attainable region and any counterfactual method-selection discrepancy was recorded without changing the frozen method.
\item[$\square$] MP-EB designs, if selected, do not treat $G$ as reducing its finite-sample range term; WSR-PP $H$/$G$ trade-off is simulation-derived.
\item[$\square$] If no design passes, feasibility-and-measurement-pilot status is frozen without weakening scientific requirements.
\end{itemize}

\subsection{Execution, provenance, and reporting}
\begin{itemize}
\item[$\square$] Final $H_p$, $G$, superblock bank selection, complete generation manifest, and root hash are frozen before confirmation.
\item[$\square$] Exactly one terminal observed/missing record exists per planned generation ID; attempts are lineage only.
\item[$\square$] Failed draws replay only on certified interchangeable devices with identical prompt bytes/seed/model/runtime.
\item[$\square$] Primary analysis consumes no post-treatment fields as covariates, weights, filters, or denominators.
\item[$\square$] $\widehat\Theta_{\rm crit}$ is labelled selected/descriptive and displayed adjacent to $L_\Theta$ and the winner-selection warning.
\item[$\square$] Separate-family multiplicity limitations and any technical-integrity warning are stated in the main result summary.
\item[$\square$] Terminal-falsification artifacts were frozen/sequestered before confirmation and opened only after locked confirmatory inference.
\end{itemize}

\section{Technical appendices}

\subsection*{Appendix A. Dated methodological and implementation evidence}
\addcontentsline{toc}{subsection}{Appendix A. Dated methodological and implementation evidence}
The MP-EB candidate uses the bounded empirical-Bernstein theorem and finite-class union logic of Maurer and Pontil \href{https://arxiv.org/abs/0907.3740}{(Maurer \& Pontil, 2009)}. The WSR-PP candidate uses the bounded conditional-mean betting/martingale construction and predictable plug-in stake design of Waudby-Smith and Ramdas \href{https://doi.org/10.1093/jrsssb/qkad009}{(Waudby-Smith \& Ramdas, 2024)}. Their sampling-without-replacement section concerns a distinct fixed finite-population model; v19 instead changes the confirmatory superblock law to independent with-replacement proposals so the ordinary conditional-mean construction applies directly. Construction bootstrap rationale follows whole-cluster resampling principles \href{https://cameron.econ.ucdavis.edu/research/Cameron_Miller_JHR_2015.pdf}{(Cameron \& Miller, 2015)}. Boundary-aware binomial calibration follows Brown, Cai, and DasGupta \href{https://www.stat.purdue.edu/~dasgupta/publications/tr99-19.pdf}{(2001)}. End-to-end statistical-method simulation follows the explicit design/evaluation discipline of Morris, White, and Crowther \href{https://doi.org/10.1002/sim.8086}{(2019)}. The selected-maximum warning is consistent with winner-selection inference concerns \href{https://doi.org/10.1093/qje/qjad043}{(Andrews, Kitagawa, \& McCloskey, 2024)}.

Instruction-following verifier precedents include IFEval \href{https://arxiv.org/abs/2311.07911}{(Zhou et al., 2023)} and ComplexBench \href{https://arxiv.org/abs/2407.03978}{(Wen et al., 2024)}. Wording sensitivity is supported by prompt/paraphrase robustness work \href{https://arxiv.org/abs/2406.19898}{(Wahle et al., 2024)}; \href{https://arxiv.org/abs/2306.04528}{(Zhu et al., 2023)}. Nucleus sampling follows Holtzman et al. \href{https://arxiv.org/abs/1904.09751}{(2020)}. Excluding realised reasoning use and final-answer length from primary adjustment follows general post-treatment-variable reasoning \href{https://doi.org/10.1111/ajps.12357}{(Montgomery, Nyhan, \& Torres, 2018)}; \href{https://doi.org/10.2307/2981697}{(Rosenbaum, 1984)}.

Implementation-dependent facts are dated 9 August 2026 and require re-verification immediately before freeze. PyTorch documents limits to reproducibility across releases/platforms and deterministic-operation controls \href{https://docs.pytorch.org/docs/stable/notes/randomness.html}{(PyTorch, 2026a)}; \href{https://docs.pytorch.org/docs/stable/generated/torch.use_deterministic_algorithms.html}{(PyTorch, 2026b)}. Hugging Face documents generation controls and chat-template tokenisation \href{https://huggingface.co/docs/transformers/main_classes/text_generation}{(Hugging Face, 2026a)}; \href{https://huggingface.co/docs/transformers/main/chat_templating}{(Hugging Face, 2026b)}. NVIDIA documents CUDA-library reproducibility conditions for a fixed toolkit/hardware context \href{https://docs.nvidia.com/cuda/cublas/index.html}{(NVIDIA, 2026)}. Lambda documents On-Demand Cloud instances, persistent filesystems, current instance-type discovery, and billing \href{https://docs.lambda.ai/public-cloud/on-demand/}{(Lambda, 2026a)}; \href{https://docs.lambda.ai/public-cloud/filesystems/}{(Lambda, 2026b)}; \href{https://docs.lambda.ai/public-cloud/cloud-api/}{(Lambda, 2026c)}; \href{https://docs.lambda.ai/public-cloud/billing/}{(Lambda, 2026d)}. vLLM documents reproducibility limitations and batch-invariance controls \href{https://docs.vllm.ai/en/latest/usage/reproducibility.html}{(vLLM Project, 2026a)}; \href{https://docs.vllm.ai/en/stable/features/batch_invariance/}{(vLLM Project, 2026b)}. Z3 documents distinct \texttt{sat}, \texttt{unsat}, and \texttt{unknown} solver results \href{https://z3prover.github.io/api/html/js/interfaces/Solver.html}{(Z3 Project, 2026)}. ROCm remains a valid local-development candidate and any support statement is reverified against the then-current official compatibility matrix \href{https://rocm.docs.amd.com/en/latest/compatibility/compatibility-matrix.html}{(AMD, 2026)}.

\subsection*{Appendix B. Reference construction in executable detail}
\addcontentsline{toc}{subsection}{Appendix B. Reference construction in executable detail}
For each pair: author 24 candidates; apply hard logical, schema, scope, satisfiability, boundary, and verifier gates; run the fixed coupling certification; certify one canonical plus 11 semantic paraphrases; estimate wording instability; require at least 12 hard-eligible candidates, otherwise add exactly one blind tranche of 12 and repeat identical hard gates; stop with \REFPOOL{} if fewer than 12 of 36 survive; collect matching coordinates on complete construction clusters $U_i$; run 10,000 cluster bootstraps; apply simultaneous coordinate calipers and the joint certificate; rank only already admissible candidates; admit exactly four; run disjoint independent validation. Every candidate has immutable IDs and rejection reason codes. No candidate can be reinstated because a later-stage outcome is convenient.

The coupling thresholds are $U_A\le0.02$, $U_R\le0.02$, and $U_\kappa\le0.015$ under the Stage~1 floor $\epsilon\ge0.06$. The initial probe bank contains 2,000 base outputs. If the frozen precision branch requires expansion until the zero-error simultaneous UCB is at most $\epsilon/8$, exact $\epsilon$ is frozen first and the bank expands only to the smallest qualifying multiple of 500. Wording admission requires $U_{|W_R-W_A|}\le0.02$ under the same SESOI floor. Any matching/caliper constant using exact $\epsilon$ is calculated only after exact $\epsilon$ has been frozen.

\subsection*{Appendix C. Confirmatory superblock population and sampling}
\addcontentsline{toc}{subsection}{Appendix C. Confirmatory superblock population and sampling}
For each pair $p$ and accepted draw index $i$, a domain-separated proposal stream independently draws: one task with replacement from $\mathcal T_p$; one background identity with replacement from each of the three frozen strata for the maximum load; all family-specific parameter realisations; and frozen construction permutations. The entire proposal is accepted only if focal and all four reference maximum-load compositions pass the complete pre-confirmatory acceptance conjunction. Otherwise the entire proposal is discarded and the same indexed stream advances to a fresh independent proposal. The first passing proposal is accepted.

The probability model is the fixed proposal law conditioned on acceptance, so accepted indices are identically distributed and independent across $i$. The implementation seed is a reproducible draw from that law rather than the rationale for independence. Repeated task/background contexts are legal and are distinguished by accepted draw index in \texttt{superblock\_id}. At $k=2$, $S=\varnothing$ and all other paired scientific fields come from the same parent superblock. Occupancy and duplicate-context frequencies are reported descriptively.

\subsection*{Appendix D. Early and final scientific-runtime validation}
\addcontentsline{toc}{subsection}{Appendix D. Early and final scientific-runtime validation}
The reference scientific architecture is PyTorch + Hugging Face Transformers + the frozen accelerator backend. The tests below are run as early as technically possible at Stage~2A/3. Every test is repeated at Stage~13 final re-certification on the selected confirmatory pool. A later behaviourally relevant runtime change invalidates the dependent downstream artifacts and requires rerunning their pre-confirmatory stages.

Mandatory tests are:
\begin{enumerate}
\item at least 100 tokenizer/chat-template fixtures with exact token-ID agreement with frozen expected vectors;
\item reasoning-cap boundary fixtures at $b-1,b,b+1$, including multi-token delimiters, immediate delimiter, no delimiter, malformed transition, and exactly-at-cap cases;
\item final-answer boundary fixtures at $V-1,V,V+1$, immediate EOS, no EOS, Unicode, and control-token cases;
\item 100,000 property-based parser/controller trajectories satisfying $N_{\rm reason}\le b$ and $N_{\rm final}\le V$ with zero invariant violations;
\item at least 1,000 same-device prompt-seed replays with exact generated token IDs, output bytes, controller flags, extracted final answer, and verifier verdict;
\item at least 1,000 cross-device replays spanning every intended GPU in the interchangeability pool with the same exact equality criterion; one mismatch fails the class;
\item at least 100,000 legal scientific IDs with deterministic seed reproduction and zero observed collisions;
\item at least 500 maximum-context runs using the complete frozen prompt plus $b_{\max}+V+$ control-token allowance with zero OOM and no silent truncation;
\item explicit dtype and attention-backend assertions with known fixtures and no silent fallback;
\item 1,000 archived-output verifier replays with byte-identical extracted answers and verdicts;
\item deliberate mutation of every critical manifest field, each required to fail closed before generation;
\item intentional worker/device death demonstrating replay of the identical \texttt{generation\_id}, prompt bytes, seed, model, and runtime on another certified device;
\item durable-persistence interruption tests proving a generation is not marked complete before raw output bytes and attempt metadata are safely committed;
\item synthetic permanent-missingness fixtures for every constituent observation type, with primitive and secondary intervals checked against brute-force enumeration of binary completions;
\item all-zero, all-one, constant interior, and all-interval inference fixtures for both MP-EB and WSR-PP, with finite support-valid output and no standard-error failure branch;
\item \textbf{validity-critical WSR-PP truncation-transition fixtures} against an independently implemented reference calculator. The frozen fixture bank must include: (a) $\widetilde\lambda_t<c/m$; (b) $\widetilde\lambda_t>c/m$; (c) equality and machine-precision neighbourhoods around $\widetilde\lambda_t=c/m$; (d) sequential transitions into and out of the truncation-bound regime; (e) exact resulting stake values, log-capital increments, total capital, rejected-mean set, and final conservative-bisection root; (f) boundary candidate means $m=0$, values arbitrarily close to $0$, values near $1$, and $m=1$; and (g) design-relevant $H$, variance, and $\delta$ configurations in which truncation is expected to bind. Production/reference disagreement on any branch identity or listed numerical output is an inference-validity failure and makes WSR-PP ineligible at Stage~9;
\item brute-force small-family verification of the complete simultaneous allocation, $L_s=\min_rL_r$, $L_\Theta=\max_sL_s$, and selected-cell tie rule;
\item null stress simulations under the frozen scenario suite demonstrating the preregistered false-positive acceptance condition;
\item exact hardware-benchmark replay sufficient to establish the selected Lambda configuration's valid-generation throughput, peak memory, failure rate, and cost metric from archived logs.
\end{enumerate}

Optional vLLM substitution requires every test above plus exactly 10,000 paired prompt-seed generations against the frozen Transformers reference. Cases span all pair types, arm types, parameter extremes, candidate/retained budgets, both orders, maximum context, and controller boundaries. All generated token IDs, persisted output bytes, controller flags, final-answer bytes, and verifier outputs must match exactly. One mismatch leaves Transformers as the scientific runtime. There is no fallback to a looser distributional-equivalence claim.

\subsection*{Appendix E. Technical missingness, retries, and interval propagation}
\addcontentsline{toc}{subsection}{Appendix E. Technical missingness, retries, and interval propagation}
Each \texttt{generation\_id} has one immutable scientific seed and one or more \texttt{attempt\_id}s. An attempt may execute on any certified-interchangeable GPU while preserving every scientific input. Once valid model output bytes are durably persisted, the generation is behaviourally observed even if the final answer is empty, malformed, refused, or cap-truncated. A deterministic verifier operational failure is repaired by rerunning verification on stored bytes. Permanent computation failure before persistence yields the exact binary-outcome support interval $[0,1]$ for that generation after the retry policy is exhausted.

For any linear superblock contrast $Q=\sum_c s_c\bar Y_c$ with known signs $s_c\in\{-1,+1\}$, compute its exact interval by selecting the lower endpoint for positive terms and upper endpoint for negative terms in the lower bound, reversing choices for the upper bound. Primitive DCI uses $[-1,1]$. Differences between two primitives use the explicit $[-2,2]$ formulas in Section~\ref{sec:secondarysupports}. The analysis ledger stores all constituent intervals, contribution intervals, point summaries, selected-method inputs, and missingness reasons.

\subsection*{Appendix F. Operational call-count and storage formulas}
\addcontentsline{toc}{subsection}{Appendix F. Operational call-count and storage formulas}
For pair $p$,
\[
N_{{\rm conf},p}=H_p|\Bret||\mathcal K||\mathcal O|(1+J)G=20H_p|\Bret|G.
\]
Across three pairs,
\[
N_{\rm conf}=20|\Bret|G\sum_{p=1}^{3}H_p.
\]
These are logical generations, not physical attempts. Horizontal GPU count changes concurrency only. Every logical generation has exactly one terminal ledger state and may have multiple attempt records. Persistent storage capacity is provisioned from the complete manifest plus a frozen safety factor before launch and must accommodate raw model bytes, controller metadata, attempt lineage, verifier records, and hashes without relying on ephemeral instance disks.

\subsection*{Appendix G. Inference-method and power-simulation schema}
\addcontentsline{toc}{subsection}{Appendix G. Inference-method and power-simulation schema}
Each scenario stores a stable ID; simulation role (method comparison or authoritative design); selected/candidate inference method; baseline target-failure vector; true primitive/reference effects; bottleneck identity; between-superblock distribution; shared focal/reference dependence; reference heterogeneity; decoder replicate model; retained budget/load/order family; boundary-contamination mechanism; technical-missingness rate and dependence mechanism; candidate $\mathbf H,G$; family cardinalities; exact supports; expected-token/runtime/cost model; and seed. The method-comparison schema additionally stores $\widetilde H_{p,\max}^{\rm plan}$, every source ceiling in $\mathcal C_p^{\rm pre}$, and the freeze hash proving those values predate comparison output. The authoritative schema stores realised $H_{p,\max}$ and the overlay of that region on the archived Stage~9 candidate grid. Every simulated experiment runs the production interval construction, selected candidate confidence code, separate-family multiplicity, and max--min logic. No analytic approximation replaces the production path for design acceptance.

For MP-EB, outputs include variance and finite-range penalty decomposition and explicitly show that only $H$ enters the range penalty. For WSR-PP, outputs include capital/root-finding diagnostics, empirical $H$/$G$ trade-offs, the fraction of stake updates that are truncation-bound, and transition counts into and out of that regime. Method comparison archives all candidates, including failed validity status; authoritative simulation runs only the frozen selected method. After Stage~10, the schema also records the method that would win if the frozen Stage~9 selection criterion were reapplied to archived candidate results restricted to $H_p\le H_{p,\max}$, together with an explicit discrepancy flag relative to the prospectively frozen method. This field is diagnostic and cannot trigger method re-selection or new simulation.

\subsection*{Appendix H. Required result outputs and figures}
\addcontentsline{toc}{subsection}{Appendix H. Required result outputs and figures}
The locked report includes: protocol/construction/runtime flow; exact selected inference method and pre-confirmatory selection report; frozen pre-Stage-10 $\widetilde H_{p,\max}^{\rm plan}$ values and their source ceilings; realised $H_{p,\max}$ overlays on archived candidate-method curves; the counterfactual restricted-attainability method winner and a prominent discrepancy statement if it differs from the method actually frozen at Stage~9; full four-reference DCI surfaces across budget/load/order; least-favourable-reference surfaces with lower bounds; $\widehat\Theta_{\rm crit}$, $s^\ast$, adjacent $L_\Theta$, and selection warning; separately protected two-sided primitive, compute, load, and order intervals with explicit familywise-guarantee wording; reference-construction audit; accepted/rejected superblock sampling and duplicate-context diagnostics; satisfiability evidential hierarchy and witness audit; verifier reliability; assigned budget versus realised reasoning use; cap/truncation/retry/technical-missingness diagnostics, including WSR-PP truncation-bound frequency if WSR-PP was a valid Stage~9 candidate; missingness interval-span and prespecified sensitivity summaries; Lambda hardware benchmark and interchangeability certification; runtime manifest and replay results; prospective method/design simulation curves; all result/status codes; and terminal-falsification result after unsequestering. Machine-readable tables are authoritative over plots.

\subsection*{Appendix I. Inference ledger and invariants}
\addcontentsline{toc}{subsection}{Appendix I. Inference ledger and invariants}
Before analysis, mechanically verify: one terminal state per planned generation ID; monotonically indexed attempts; every successful generation linked to durably stored output bytes; every permanent missing value linked to exhausted retry lineage; no attempt executed on an uncertified device/runtime; focal/reference pretreatment equality except the added-slot occupant and its own fields; common parent superblock across paired loads/orders; all four references at every required cell; paired parameters; budget/order membership in frozen sets; invariant model/tokenizer/template/controller/decoder/runtime/verifier/inference hashes; no primary formula consumes post-treatment/diagnostic fields; $H_r=H_p$ for each primitive; primitive support exactly $[-1,1]$; secondary-confirmatory support exactly $[-2,2]$; primary family cardinality $48|\Bret|$; if WSR-PP is selected, primary memberwise $\delta$ exactly $0.05/(48|\Bret|)$; compute family $48$; load and order families $24|\Bret|$ each; separately controlled familywise alpha recorded for every family; selected inference method identical to the pre-confirmatory method-selection artifact; deterministic reproduction of all $L_r$, upper bounds where relevant, $L_s$, $L_\Theta$, $s^\ast$, and status codes. Violations emit the appropriate fail-closed state rather than being silently repaired.

\subsection*{Appendix J. Construction, runtime, inference, and freeze artifact manifest}
\addcontentsline{toc}{subsection}{Appendix J. Construction, runtime, inference, and freeze artifact manifest}
The freeze bundle stores: constraint registries and parameter grids; material split manifests; verifier code/audits; candidate references and hard-gate records; coupling bank/results; semantic wording certifications; wording data/bootstrap outputs; matching calibration data/seeds/constants; admitted panels and independent validation; checkpoint-selection record; exact model/tokenizer/template files and hashes; Stage~2A validation; Lambda candidate list and hardware benchmark; exact accelerator/GPU/runtime manifest; interchangeability-device roster and cross-device replay results; budget manipulation and $V$ artifacts; exact SESOI, signed scientific rationale, and timestamp of first dependency; inference-candidate theorem mappings, validity-critical WSR-PP truncation fixtures, $\widetilde H_{p,\max}^{\rm plan}$ source ceilings, 10,000-experiment comparison, selected method/configuration/code hashes; every superblock proposal and satisfiability/witness evidence; attainable banks and $H_{p,\max}$; realised-attainability counterfactual method-selection diagnostic; authoritative power/design simulations; final $H_p,G$; complete scientific-draw manifest and seeds; persistent-storage layout; final runtime re-certification; analysis code; preregistration source/PDF; remaining-pre-freeze closure report; and sequestered terminal-falsification plan. A root SHA-256 manifest is signed before the first confirmatory output.

\subsection*{Appendix K. References}
\addcontentsline{toc}{subsection}{Appendix K. References}
\begin{hangparas}{1.5em}{1}
Advanced Micro Devices, Inc. (2026). \textit{ROCm compatibility matrix}. Retrieved August 9, 2026, from \url{https://rocm.docs.amd.com/en/latest/compatibility/compatibility-matrix.html}

Andrews, I., Kitagawa, T., \& McCloskey, A. (2024). Inference on winners. \textit{The Quarterly Journal of Economics, 139}(1), 305--358. \url{https://doi.org/10.1093/qje/qjad043}

Brown, L. D., Cai, T. T., \& DasGupta, A. (2001). Interval estimation for a binomial proportion. \textit{Statistical Science, 16}(2), 101--133. \url{https://www.stat.purdue.edu/~dasgupta/publications/tr99-19.pdf}

Cameron, A. C., \& Miller, D. L. (2015). A practitioner's guide to cluster-robust inference. \textit{Journal of Human Resources, 50}(2), 317--372. \url{https://cameron.econ.ucdavis.edu/research/Cameron_Miller_JHR_2015.pdf}

Holtzman, A., Buys, J., Du, L., Forbes, M., \& Choi, Y. (2020). The curious case of neural text degeneration. \textit{International Conference on Learning Representations}. \url{https://arxiv.org/abs/1904.09751}

Hugging Face. (2026a). \textit{Generation}. Transformers documentation. Retrieved August 9, 2026, from \url{https://huggingface.co/docs/transformers/main_classes/text_generation}

Hugging Face. (2026b). \textit{Chat templates}. Transformers documentation. Retrieved August 9, 2026, from \url{https://huggingface.co/docs/transformers/main/chat_templating}

Lakens, D. (2017). Equivalence tests: A practical primer for $t$ tests, correlations, and meta-analyses. \textit{Social Psychological and Personality Science, 8}(4), 355--362. \url{https://doi.org/10.1177/1948550617697177}

Lambda. (2026a). \textit{On-Demand Cloud overview}. Retrieved August 9, 2026, from \url{https://docs.lambda.ai/public-cloud/on-demand/}

Lambda. (2026b). \textit{Filesystems}. Retrieved August 9, 2026, from \url{https://docs.lambda.ai/public-cloud/filesystems/}

Lambda. (2026c). \textit{Lambda Cloud API: Instance types and instance operations}. Retrieved August 9, 2026, from \url{https://docs.lambda.ai/public-cloud/cloud-api/}

Lambda. (2026d). \textit{Billing overview}. Retrieved August 9, 2026, from \url{https://docs.lambda.ai/public-cloud/billing/}

Lampropoulos, L., Hicks, M., \& Pierce, B. C. (2019). Coverage guided, property based testing. \textit{Proceedings of the ACM on Programming Languages, 3}(OOPSLA). \url{https://doi.org/10.1145/3360607}

Maurer, A., \& Pontil, M. (2009). Empirical Bernstein bounds and sample variance penalization. \textit{Proceedings of the 22nd Conference on Learning Theory}. \url{https://arxiv.org/abs/0907.3740}

Montgomery, J. M., Nyhan, B., \& Torres, M. (2018). How conditioning on posttreatment variables can ruin your experiment and what to do about it. \textit{American Journal of Political Science, 62}(3), 760--775. \url{https://doi.org/10.1111/ajps.12357}

Morris, T. P., White, I. R., \& Crowther, M. J. (2019). Using simulation studies to evaluate statistical methods. \textit{Statistics in Medicine, 38}, 2074--2102. \url{https://doi.org/10.1002/sim.8086}

NVIDIA. (2026). \textit{cuBLAS documentation: Results reproducibility}. Retrieved August 9, 2026, from \url{https://docs.nvidia.com/cuda/cublas/index.html}

PyTorch. (2026a). \textit{Reproducibility}. Retrieved August 9, 2026, from \url{https://docs.pytorch.org/docs/stable/notes/randomness.html}

PyTorch. (2026b). \textit{torch.use\_deterministic\_algorithms}. Retrieved August 9, 2026, from \url{https://docs.pytorch.org/docs/stable/generated/torch.use_deterministic_algorithms.html}

Rosenbaum, P. R. (1984). The consequences of adjustment for a concomitant variable that has been affected by the treatment. \textit{Journal of the Royal Statistical Society: Series A, 147}(5), 656--666. \url{https://doi.org/10.2307/2981697}

VanderWeele, T. J. (2016). Mediation analysis: A practitioner's guide. \textit{Annual Review of Public Health, 37}, 17--32. \url{https://doi.org/10.1146/annurev-publhealth-032315-021402}

Vasileva, M. I. (2026). \textit{Large language models can follow instructions, but not many at once: Phase transitions in compositional constraint satisfaction}. OpenReview. \url{https://openreview.net/forum?id=a6r2ECEVls}

vLLM Project. (2026a). \textit{Reproducibility}. Retrieved August 9, 2026, from \url{https://docs.vllm.ai/en/latest/usage/reproducibility.html}

vLLM Project. (2026b). \textit{Batch invariance}. Retrieved August 9, 2026, from \url{https://docs.vllm.ai/en/stable/features/batch_invariance/}

Wahle, J. P., Ruas, T., Xu, Y., \& Gipp, B. (2024). Paraphrase types elicit prompt engineering capabilities. \textit{Proceedings of EMNLP 2024}. \url{https://arxiv.org/abs/2406.19898}

Waudby-Smith, I., \& Ramdas, A. (2024). Estimating means of bounded random variables by betting. \textit{Journal of the Royal Statistical Society: Series B (Statistical Methodology), 86}(1), 1--27. \url{https://doi.org/10.1093/jrsssb/qkad009}

Wen, B., et al. (2024). Benchmarking complex instruction-following with multiple constraints composition. \url{https://arxiv.org/abs/2407.03978}

Z3 Project. (2026). \textit{Z3 solver API: satisfiability results}. Retrieved August 9, 2026, from \url{https://z3prover.github.io/api/html/js/interfaces/Solver.html}

Zhou, J., Lu, T., Mishra, S., Brahma, S., Basu, S., Luan, Y., Zhou, D., \& Hou, L. (2023). Instruction-following evaluation for large language models. \url{https://arxiv.org/abs/2311.07911}

Zhu, K., et al. (2023). PromptRobust: Towards evaluating the robustness of large language models on adversarial prompts. \url{https://arxiv.org/abs/2306.04528}
\end{hangparas}

\clearpage\section*{Final comprehension check}
\addcontentsline{toc}{section}{Final comprehension check}
A technically literate reader should be able to determine from this protocol alone: the replacement estimand; the four-reference admission logic; the pre-frozen population after boundary/satisfiability construction; the stochastic decoder and hard-budget controller; the distinction between a retry and a scientific redraw; interval handling of permanent technical failures; the two pre-confirmatory finite-sample inference candidates and the validity-first rule selecting one; why zero variance remains estimable; how simultaneous primitive bounds imply $L_\Theta\le\Theta_{\rm crit}$; how paired order effects are tested; how $H_p$ and $G$ are selected; which runtime tests must pass; which structural failures can truly create \INFERNE{}; and which numerical/artifact choices remain pre-freeze.

\vfill
\begin{center}
\textit{End of v19-pilot DRAFT preregistration. Protocol remains pre-freeze until every freeze-blocking item in Section~\ref{sec:register} is resolved and the immutable freeze stage is executed.}
\end{center}

\end{document}
